\documentclass[11pt]{article}

% Shared notes settings for Applied Machine Learning lectures (UPES).
% Keep this file free of \documentclass to allow reuse via \input.

\usepackage[utf8]{inputenc}
\usepackage[T1]{fontenc}
\usepackage{lmodern}          % scalable T1 fonts; without this pdflatex falls back
\input{glyphtounicode}        % to bitmapped EC Type 3 and the PDFs are neither
\pdfgentounicode=1            % searchable nor screen-reader friendly
\usepackage[a4paper,margin=1in]{geometry}
\usepackage{hyperref}
\usepackage{xurl}
\usepackage{booktabs}
\usepackage{amsmath}
\usepackage{amssymb}
\usepackage{listings}
\usepackage{xcolor}
\usepackage{enumitem}
\usepackage{parskip}

% Code listings (Python-oriented for this subject).
\lstset{
  language=Python,
  basicstyle=\ttfamily\small,
  keywordstyle=\color{blue},
  commentstyle=\color{gray},
  stringstyle=\color{teal},
  showstringspaces=false,
  columns=fullflexible,
  frame=single,
  framerule=0.2pt,
  breaklines=true
}

\setlist[itemize]{topsep=0.3em,itemsep=0.2em}
\setlist[enumerate]{topsep=0.3em,itemsep=0.2em}

% Simple per-section source line.
\newcommand{\notesources}[1]{\par\noindent{\small\color{gray}\textit{Sources:} \sloppy #1\par}}

% Unit-wise source strings for this subject (used by slides + notes).
% Path is relative to the lecture `latex/` folder (the compilation CWD).
\IfFileExists{../../../_shared/aml-sources.tex}{%
  % Source strings for Applied Machine Learning (CSAI2017P) — used by slides + notes.
% Prescribed texts and reference books are taken from the course syllabus.
% Support texts are the author-hosted open-access editions:
%   statlearning.com            (ISL with Python)
%   probml.github.io/pml-book   (Murphy, Probabilistic Machine Learning)
%   d2l.ai                      (Dive into Deep Learning)
%   mml-book.github.io          (Mathematics for Machine Learning)

\newcommand{\AMLRepoURL}{https://github.com/tali7c/Applied-Machine-Learning}

% --- prescribed -------------------------------------------------------------
\newcommand{\AMLPrescribed}{%
M\"uller and Guido, \textit{Introduction to Machine Learning with Python}, O'Reilly, 2016.;%
Bishop, \textit{Pattern Recognition and Machine Learning}, Springer, 2006.;%
Murphy, \textit{Machine Learning: A Probabilistic Perspective}, MIT Press, 2012.;%
Han, Kamber and Pei, \textit{Data Mining: Concepts and Techniques} (3e), Morgan Kaufmann, 2012.%
}

% --- unit-wise ---------------------------------------------------------------
\newcommand{\AMLUnitISources}{%
M\"uller and Guido, \textit{Introduction to Machine Learning with Python}, Ch. 1.;%
Bishop, \textit{Pattern Recognition and Machine Learning}, Ch. 1.;%
Murphy, \textit{Probabilistic Machine Learning: An Introduction}, MIT Press, 2022, Ch. 1.;%
Mitchell, \textit{Machine Learning}, McGraw Hill, 1997, Ch. 1.;%
Scikit-learn documentation, accessed 2026-08-04.%
}

\newcommand{\AMLUnitIISources}{%
Bishop, \textit{Pattern Recognition and Machine Learning}, \S1.5, \S4.3, \S7.1.;%
Zhang, Lipton, Li and Smola, \textit{Dive into Deep Learning}, \S3.4.;%
Shalev-Shwartz and Ben-David, \textit{Understanding Machine Learning}, Ch. 15.;%
Murphy, \textit{Probabilistic Machine Learning: An Introduction}, Ch. 5.%
}

\newcommand{\AMLUnitIIISources}{%
Zhang, Lipton, Li and Smola, \textit{Dive into Deep Learning}, Ch. 12 (Optimization Algorithms).;%
Deisenroth, Faisal and Ong, \textit{Mathematics for Machine Learning}, Ch. 7.;%
Kingma and Ba, ``Adam: A Method for Stochastic Optimization,'' ICLR 2015.;%
Ruder, ``An overview of gradient descent optimization algorithms,'' arXiv:1609.04747, 2016.%
}

\newcommand{\AMLUnitIVSources}{%
Han, Kamber and Pei, \textit{Data Mining: Concepts and Techniques} (3e), Ch. 3.;%
M\"uller and Guido, \textit{Introduction to Machine Learning with Python}, Ch. 3--4.;%
James, Witten, Hastie, Tibshirani and Taylor, \textit{An Introduction to Statistical Learning with Python}, Ch. 5--6, 12.;%
Scikit-learn documentation (preprocessing, model selection), accessed 2026-08-04.%
}

\newcommand{\AMLUnitVSources}{%
James et al., \textit{An Introduction to Statistical Learning with Python}, Ch. 3--6.;%
Bishop, \textit{Pattern Recognition and Machine Learning}, Ch. 3.;%
Deisenroth, Faisal and Ong, \textit{Mathematics for Machine Learning}, Ch. 9.;%
M\"uller and Guido, \textit{Introduction to Machine Learning with Python}, \S2.3.%
}

\newcommand{\AMLUnitVISources}{%
Han, Kamber and Pei, \textit{Data Mining: Concepts and Techniques} (3e), Ch. 8--9.;%
James et al., \textit{An Introduction to Statistical Learning with Python}, Ch. 8--9.;%
Bishop, \textit{Pattern Recognition and Machine Learning}, Ch. 4--5, 7, 14.;%
Zhang et al., \textit{Dive into Deep Learning}, Ch. 5.%
}

\newcommand{\AMLUnitVIISources}{%
Han, Kamber and Pei, \textit{Data Mining: Concepts and Techniques} (3e), Ch. 10.;%
James et al., \textit{An Introduction to Statistical Learning with Python}, Ch. 12.;%
M\"uller and Guido, \textit{Introduction to Machine Learning with Python}, \S3.5.;%
Ester, Kriegel, Sander and Xu, ``A density-based algorithm for discovering clusters (DBSCAN),'' KDD 1996.%
}

\newcommand{\AMLLabSources}{%
M\"uller and Guido, \textit{Introduction to Machine Learning with Python}, companion notebooks.;%
McKinney, \textit{Python for Data Analysis} (3e), O'Reilly, 2022.;%
Scikit-learn user guide, accessed 2026-08-04.;%
Matplotlib and Seaborn documentation, accessed 2026-08-04.%
}
%
}{%
  % Faculty copies live under `private/` so their compile CWD is different.
  \IfFileExists{../../../../public/_shared/aml-sources.tex}{%
    % Source strings for Applied Machine Learning (CSAI2017P) — used by slides + notes.
% Prescribed texts and reference books are taken from the course syllabus.
% Support texts are the author-hosted open-access editions:
%   statlearning.com            (ISL with Python)
%   probml.github.io/pml-book   (Murphy, Probabilistic Machine Learning)
%   d2l.ai                      (Dive into Deep Learning)
%   mml-book.github.io          (Mathematics for Machine Learning)

\newcommand{\AMLRepoURL}{https://github.com/tali7c/Applied-Machine-Learning}

% --- prescribed -------------------------------------------------------------
\newcommand{\AMLPrescribed}{%
M\"uller and Guido, \textit{Introduction to Machine Learning with Python}, O'Reilly, 2016.;%
Bishop, \textit{Pattern Recognition and Machine Learning}, Springer, 2006.;%
Murphy, \textit{Machine Learning: A Probabilistic Perspective}, MIT Press, 2012.;%
Han, Kamber and Pei, \textit{Data Mining: Concepts and Techniques} (3e), Morgan Kaufmann, 2012.%
}

% --- unit-wise ---------------------------------------------------------------
\newcommand{\AMLUnitISources}{%
M\"uller and Guido, \textit{Introduction to Machine Learning with Python}, Ch. 1.;%
Bishop, \textit{Pattern Recognition and Machine Learning}, Ch. 1.;%
Murphy, \textit{Probabilistic Machine Learning: An Introduction}, MIT Press, 2022, Ch. 1.;%
Mitchell, \textit{Machine Learning}, McGraw Hill, 1997, Ch. 1.;%
Scikit-learn documentation, accessed 2026-08-04.%
}

\newcommand{\AMLUnitIISources}{%
Bishop, \textit{Pattern Recognition and Machine Learning}, \S1.5, \S4.3, \S7.1.;%
Zhang, Lipton, Li and Smola, \textit{Dive into Deep Learning}, \S3.4.;%
Shalev-Shwartz and Ben-David, \textit{Understanding Machine Learning}, Ch. 15.;%
Murphy, \textit{Probabilistic Machine Learning: An Introduction}, Ch. 5.%
}

\newcommand{\AMLUnitIIISources}{%
Zhang, Lipton, Li and Smola, \textit{Dive into Deep Learning}, Ch. 12 (Optimization Algorithms).;%
Deisenroth, Faisal and Ong, \textit{Mathematics for Machine Learning}, Ch. 7.;%
Kingma and Ba, ``Adam: A Method for Stochastic Optimization,'' ICLR 2015.;%
Ruder, ``An overview of gradient descent optimization algorithms,'' arXiv:1609.04747, 2016.%
}

\newcommand{\AMLUnitIVSources}{%
Han, Kamber and Pei, \textit{Data Mining: Concepts and Techniques} (3e), Ch. 3.;%
M\"uller and Guido, \textit{Introduction to Machine Learning with Python}, Ch. 3--4.;%
James, Witten, Hastie, Tibshirani and Taylor, \textit{An Introduction to Statistical Learning with Python}, Ch. 5--6, 12.;%
Scikit-learn documentation (preprocessing, model selection), accessed 2026-08-04.%
}

\newcommand{\AMLUnitVSources}{%
James et al., \textit{An Introduction to Statistical Learning with Python}, Ch. 3--6.;%
Bishop, \textit{Pattern Recognition and Machine Learning}, Ch. 3.;%
Deisenroth, Faisal and Ong, \textit{Mathematics for Machine Learning}, Ch. 9.;%
M\"uller and Guido, \textit{Introduction to Machine Learning with Python}, \S2.3.%
}

\newcommand{\AMLUnitVISources}{%
Han, Kamber and Pei, \textit{Data Mining: Concepts and Techniques} (3e), Ch. 8--9.;%
James et al., \textit{An Introduction to Statistical Learning with Python}, Ch. 8--9.;%
Bishop, \textit{Pattern Recognition and Machine Learning}, Ch. 4--5, 7, 14.;%
Zhang et al., \textit{Dive into Deep Learning}, Ch. 5.%
}

\newcommand{\AMLUnitVIISources}{%
Han, Kamber and Pei, \textit{Data Mining: Concepts and Techniques} (3e), Ch. 10.;%
James et al., \textit{An Introduction to Statistical Learning with Python}, Ch. 12.;%
M\"uller and Guido, \textit{Introduction to Machine Learning with Python}, \S3.5.;%
Ester, Kriegel, Sander and Xu, ``A density-based algorithm for discovering clusters (DBSCAN),'' KDD 1996.%
}

\newcommand{\AMLLabSources}{%
M\"uller and Guido, \textit{Introduction to Machine Learning with Python}, companion notebooks.;%
McKinney, \textit{Python for Data Analysis} (3e), O'Reilly, 2022.;%
Scikit-learn user guide, accessed 2026-08-04.;%
Matplotlib and Seaborn documentation, accessed 2026-08-04.%
}
%
  }{%
    % Question papers live at `private/<family>/<instrument>/`.
    \IfFileExists{../../../public/_shared/aml-sources.tex}{%
      % Source strings for Applied Machine Learning (CSAI2017P) — used by slides + notes.
% Prescribed texts and reference books are taken from the course syllabus.
% Support texts are the author-hosted open-access editions:
%   statlearning.com            (ISL with Python)
%   probml.github.io/pml-book   (Murphy, Probabilistic Machine Learning)
%   d2l.ai                      (Dive into Deep Learning)
%   mml-book.github.io          (Mathematics for Machine Learning)

\newcommand{\AMLRepoURL}{https://github.com/tali7c/Applied-Machine-Learning}

% --- prescribed -------------------------------------------------------------
\newcommand{\AMLPrescribed}{%
M\"uller and Guido, \textit{Introduction to Machine Learning with Python}, O'Reilly, 2016.;%
Bishop, \textit{Pattern Recognition and Machine Learning}, Springer, 2006.;%
Murphy, \textit{Machine Learning: A Probabilistic Perspective}, MIT Press, 2012.;%
Han, Kamber and Pei, \textit{Data Mining: Concepts and Techniques} (3e), Morgan Kaufmann, 2012.%
}

% --- unit-wise ---------------------------------------------------------------
\newcommand{\AMLUnitISources}{%
M\"uller and Guido, \textit{Introduction to Machine Learning with Python}, Ch. 1.;%
Bishop, \textit{Pattern Recognition and Machine Learning}, Ch. 1.;%
Murphy, \textit{Probabilistic Machine Learning: An Introduction}, MIT Press, 2022, Ch. 1.;%
Mitchell, \textit{Machine Learning}, McGraw Hill, 1997, Ch. 1.;%
Scikit-learn documentation, accessed 2026-08-04.%
}

\newcommand{\AMLUnitIISources}{%
Bishop, \textit{Pattern Recognition and Machine Learning}, \S1.5, \S4.3, \S7.1.;%
Zhang, Lipton, Li and Smola, \textit{Dive into Deep Learning}, \S3.4.;%
Shalev-Shwartz and Ben-David, \textit{Understanding Machine Learning}, Ch. 15.;%
Murphy, \textit{Probabilistic Machine Learning: An Introduction}, Ch. 5.%
}

\newcommand{\AMLUnitIIISources}{%
Zhang, Lipton, Li and Smola, \textit{Dive into Deep Learning}, Ch. 12 (Optimization Algorithms).;%
Deisenroth, Faisal and Ong, \textit{Mathematics for Machine Learning}, Ch. 7.;%
Kingma and Ba, ``Adam: A Method for Stochastic Optimization,'' ICLR 2015.;%
Ruder, ``An overview of gradient descent optimization algorithms,'' arXiv:1609.04747, 2016.%
}

\newcommand{\AMLUnitIVSources}{%
Han, Kamber and Pei, \textit{Data Mining: Concepts and Techniques} (3e), Ch. 3.;%
M\"uller and Guido, \textit{Introduction to Machine Learning with Python}, Ch. 3--4.;%
James, Witten, Hastie, Tibshirani and Taylor, \textit{An Introduction to Statistical Learning with Python}, Ch. 5--6, 12.;%
Scikit-learn documentation (preprocessing, model selection), accessed 2026-08-04.%
}

\newcommand{\AMLUnitVSources}{%
James et al., \textit{An Introduction to Statistical Learning with Python}, Ch. 3--6.;%
Bishop, \textit{Pattern Recognition and Machine Learning}, Ch. 3.;%
Deisenroth, Faisal and Ong, \textit{Mathematics for Machine Learning}, Ch. 9.;%
M\"uller and Guido, \textit{Introduction to Machine Learning with Python}, \S2.3.%
}

\newcommand{\AMLUnitVISources}{%
Han, Kamber and Pei, \textit{Data Mining: Concepts and Techniques} (3e), Ch. 8--9.;%
James et al., \textit{An Introduction to Statistical Learning with Python}, Ch. 8--9.;%
Bishop, \textit{Pattern Recognition and Machine Learning}, Ch. 4--5, 7, 14.;%
Zhang et al., \textit{Dive into Deep Learning}, Ch. 5.%
}

\newcommand{\AMLUnitVIISources}{%
Han, Kamber and Pei, \textit{Data Mining: Concepts and Techniques} (3e), Ch. 10.;%
James et al., \textit{An Introduction to Statistical Learning with Python}, Ch. 12.;%
M\"uller and Guido, \textit{Introduction to Machine Learning with Python}, \S3.5.;%
Ester, Kriegel, Sander and Xu, ``A density-based algorithm for discovering clusters (DBSCAN),'' KDD 1996.%
}

\newcommand{\AMLLabSources}{%
M\"uller and Guido, \textit{Introduction to Machine Learning with Python}, companion notebooks.;%
McKinney, \textit{Python for Data Analysis} (3e), O'Reilly, 2022.;%
Scikit-learn user guide, accessed 2026-08-04.;%
Matplotlib and Seaborn documentation, accessed 2026-08-04.%
}
%
    }{%
      % Marking schemes live at `private/solutions/<family>/<id>/latex/`.
      \IfFileExists{../../../../../public/_shared/aml-sources.tex}{%
        % Source strings for Applied Machine Learning (CSAI2017P) — used by slides + notes.
% Prescribed texts and reference books are taken from the course syllabus.
% Support texts are the author-hosted open-access editions:
%   statlearning.com            (ISL with Python)
%   probml.github.io/pml-book   (Murphy, Probabilistic Machine Learning)
%   d2l.ai                      (Dive into Deep Learning)
%   mml-book.github.io          (Mathematics for Machine Learning)

\newcommand{\AMLRepoURL}{https://github.com/tali7c/Applied-Machine-Learning}

% --- prescribed -------------------------------------------------------------
\newcommand{\AMLPrescribed}{%
M\"uller and Guido, \textit{Introduction to Machine Learning with Python}, O'Reilly, 2016.;%
Bishop, \textit{Pattern Recognition and Machine Learning}, Springer, 2006.;%
Murphy, \textit{Machine Learning: A Probabilistic Perspective}, MIT Press, 2012.;%
Han, Kamber and Pei, \textit{Data Mining: Concepts and Techniques} (3e), Morgan Kaufmann, 2012.%
}

% --- unit-wise ---------------------------------------------------------------
\newcommand{\AMLUnitISources}{%
M\"uller and Guido, \textit{Introduction to Machine Learning with Python}, Ch. 1.;%
Bishop, \textit{Pattern Recognition and Machine Learning}, Ch. 1.;%
Murphy, \textit{Probabilistic Machine Learning: An Introduction}, MIT Press, 2022, Ch. 1.;%
Mitchell, \textit{Machine Learning}, McGraw Hill, 1997, Ch. 1.;%
Scikit-learn documentation, accessed 2026-08-04.%
}

\newcommand{\AMLUnitIISources}{%
Bishop, \textit{Pattern Recognition and Machine Learning}, \S1.5, \S4.3, \S7.1.;%
Zhang, Lipton, Li and Smola, \textit{Dive into Deep Learning}, \S3.4.;%
Shalev-Shwartz and Ben-David, \textit{Understanding Machine Learning}, Ch. 15.;%
Murphy, \textit{Probabilistic Machine Learning: An Introduction}, Ch. 5.%
}

\newcommand{\AMLUnitIIISources}{%
Zhang, Lipton, Li and Smola, \textit{Dive into Deep Learning}, Ch. 12 (Optimization Algorithms).;%
Deisenroth, Faisal and Ong, \textit{Mathematics for Machine Learning}, Ch. 7.;%
Kingma and Ba, ``Adam: A Method for Stochastic Optimization,'' ICLR 2015.;%
Ruder, ``An overview of gradient descent optimization algorithms,'' arXiv:1609.04747, 2016.%
}

\newcommand{\AMLUnitIVSources}{%
Han, Kamber and Pei, \textit{Data Mining: Concepts and Techniques} (3e), Ch. 3.;%
M\"uller and Guido, \textit{Introduction to Machine Learning with Python}, Ch. 3--4.;%
James, Witten, Hastie, Tibshirani and Taylor, \textit{An Introduction to Statistical Learning with Python}, Ch. 5--6, 12.;%
Scikit-learn documentation (preprocessing, model selection), accessed 2026-08-04.%
}

\newcommand{\AMLUnitVSources}{%
James et al., \textit{An Introduction to Statistical Learning with Python}, Ch. 3--6.;%
Bishop, \textit{Pattern Recognition and Machine Learning}, Ch. 3.;%
Deisenroth, Faisal and Ong, \textit{Mathematics for Machine Learning}, Ch. 9.;%
M\"uller and Guido, \textit{Introduction to Machine Learning with Python}, \S2.3.%
}

\newcommand{\AMLUnitVISources}{%
Han, Kamber and Pei, \textit{Data Mining: Concepts and Techniques} (3e), Ch. 8--9.;%
James et al., \textit{An Introduction to Statistical Learning with Python}, Ch. 8--9.;%
Bishop, \textit{Pattern Recognition and Machine Learning}, Ch. 4--5, 7, 14.;%
Zhang et al., \textit{Dive into Deep Learning}, Ch. 5.%
}

\newcommand{\AMLUnitVIISources}{%
Han, Kamber and Pei, \textit{Data Mining: Concepts and Techniques} (3e), Ch. 10.;%
James et al., \textit{An Introduction to Statistical Learning with Python}, Ch. 12.;%
M\"uller and Guido, \textit{Introduction to Machine Learning with Python}, \S3.5.;%
Ester, Kriegel, Sander and Xu, ``A density-based algorithm for discovering clusters (DBSCAN),'' KDD 1996.%
}

\newcommand{\AMLLabSources}{%
M\"uller and Guido, \textit{Introduction to Machine Learning with Python}, companion notebooks.;%
McKinney, \textit{Python for Data Analysis} (3e), O'Reilly, 2022.;%
Scikit-learn user guide, accessed 2026-08-04.;%
Matplotlib and Seaborn documentation, accessed 2026-08-04.%
}
%
      }{%
        % Source strings for Applied Machine Learning (CSAI2017P) — used by slides + notes.
% Prescribed texts and reference books are taken from the course syllabus.
% Support texts are the author-hosted open-access editions:
%   statlearning.com            (ISL with Python)
%   probml.github.io/pml-book   (Murphy, Probabilistic Machine Learning)
%   d2l.ai                      (Dive into Deep Learning)
%   mml-book.github.io          (Mathematics for Machine Learning)

\newcommand{\AMLRepoURL}{https://github.com/tali7c/Applied-Machine-Learning}

% --- prescribed -------------------------------------------------------------
\newcommand{\AMLPrescribed}{%
M\"uller and Guido, \textit{Introduction to Machine Learning with Python}, O'Reilly, 2016.;%
Bishop, \textit{Pattern Recognition and Machine Learning}, Springer, 2006.;%
Murphy, \textit{Machine Learning: A Probabilistic Perspective}, MIT Press, 2012.;%
Han, Kamber and Pei, \textit{Data Mining: Concepts and Techniques} (3e), Morgan Kaufmann, 2012.%
}

% --- unit-wise ---------------------------------------------------------------
\newcommand{\AMLUnitISources}{%
M\"uller and Guido, \textit{Introduction to Machine Learning with Python}, Ch. 1.;%
Bishop, \textit{Pattern Recognition and Machine Learning}, Ch. 1.;%
Murphy, \textit{Probabilistic Machine Learning: An Introduction}, MIT Press, 2022, Ch. 1.;%
Mitchell, \textit{Machine Learning}, McGraw Hill, 1997, Ch. 1.;%
Scikit-learn documentation, accessed 2026-08-04.%
}

\newcommand{\AMLUnitIISources}{%
Bishop, \textit{Pattern Recognition and Machine Learning}, \S1.5, \S4.3, \S7.1.;%
Zhang, Lipton, Li and Smola, \textit{Dive into Deep Learning}, \S3.4.;%
Shalev-Shwartz and Ben-David, \textit{Understanding Machine Learning}, Ch. 15.;%
Murphy, \textit{Probabilistic Machine Learning: An Introduction}, Ch. 5.%
}

\newcommand{\AMLUnitIIISources}{%
Zhang, Lipton, Li and Smola, \textit{Dive into Deep Learning}, Ch. 12 (Optimization Algorithms).;%
Deisenroth, Faisal and Ong, \textit{Mathematics for Machine Learning}, Ch. 7.;%
Kingma and Ba, ``Adam: A Method for Stochastic Optimization,'' ICLR 2015.;%
Ruder, ``An overview of gradient descent optimization algorithms,'' arXiv:1609.04747, 2016.%
}

\newcommand{\AMLUnitIVSources}{%
Han, Kamber and Pei, \textit{Data Mining: Concepts and Techniques} (3e), Ch. 3.;%
M\"uller and Guido, \textit{Introduction to Machine Learning with Python}, Ch. 3--4.;%
James, Witten, Hastie, Tibshirani and Taylor, \textit{An Introduction to Statistical Learning with Python}, Ch. 5--6, 12.;%
Scikit-learn documentation (preprocessing, model selection), accessed 2026-08-04.%
}

\newcommand{\AMLUnitVSources}{%
James et al., \textit{An Introduction to Statistical Learning with Python}, Ch. 3--6.;%
Bishop, \textit{Pattern Recognition and Machine Learning}, Ch. 3.;%
Deisenroth, Faisal and Ong, \textit{Mathematics for Machine Learning}, Ch. 9.;%
M\"uller and Guido, \textit{Introduction to Machine Learning with Python}, \S2.3.%
}

\newcommand{\AMLUnitVISources}{%
Han, Kamber and Pei, \textit{Data Mining: Concepts and Techniques} (3e), Ch. 8--9.;%
James et al., \textit{An Introduction to Statistical Learning with Python}, Ch. 8--9.;%
Bishop, \textit{Pattern Recognition and Machine Learning}, Ch. 4--5, 7, 14.;%
Zhang et al., \textit{Dive into Deep Learning}, Ch. 5.%
}

\newcommand{\AMLUnitVIISources}{%
Han, Kamber and Pei, \textit{Data Mining: Concepts and Techniques} (3e), Ch. 10.;%
James et al., \textit{An Introduction to Statistical Learning with Python}, Ch. 12.;%
M\"uller and Guido, \textit{Introduction to Machine Learning with Python}, \S3.5.;%
Ester, Kriegel, Sander and Xu, ``A density-based algorithm for discovering clusters (DBSCAN),'' KDD 1996.%
}

\newcommand{\AMLLabSources}{%
M\"uller and Guido, \textit{Introduction to Machine Learning with Python}, companion notebooks.;%
McKinney, \textit{Python for Data Analysis} (3e), O'Reilly, 2022.;%
Scikit-learn user guide, accessed 2026-08-04.;%
Matplotlib and Seaborn documentation, accessed 2026-08-04.%
}
%
      }%
    }%
  }%
}


\usepackage{array}
\usepackage{multicol}
\hypersetup{colorlinks=true, linkcolor=UPESBlueD, urlcolor=UPESBlueD}

\newcommand{\method}[1]{\begin{keyidea}[The method: #1]}
\newcommand{\ans}[1]{\subsection*{#1}}
\newcommand{\sgn}{\operatorname{sign}}
\newcommand{\IQR}{\mathrm{IQR}}
\newcommand{\MAD}{\mathrm{MAD}}
\newcommand{\unit}[1]{\par\medskip\noindent{\large\bfseries\color{UPESPurple}#1}\par\smallskip}

\newtcolorbox{concept}[1][Short-answer questions, with model answers]{breakable, enhanced,
  colback=black!2, colframe=black!45, fonttitle=\bfseries\small,
  title={#1}, boxrule=0.6pt, arc=2pt, left=6pt,right=6pt,top=4pt,bottom=4pt}

\title{Applied Machine Learning (CSAI2017P)\\
  \large Mid-Semester Study Material\\
  \normalsize The examinable topics of Lectures 1--9, worked by hand ---
  with the course's practice questions, exercises and answers}
\author{Dr.\ Tofik Ali \\ School of Computer Science, UPES Dehradun}
\date{Autumn 2026 \textbullet\ Lecture packages L1--L9 (Units I--IV and Unit V to Lecture 9)}

\begin{document}
\sloppy
\maketitle

\begin{keyidea}[{{Read this first}}]
\textbf{Syllabus.} The mid-semester paper covers Lectures 1--9:
Unit I (L1), Unit II (L2), Unit III (L3--L4), Unit IV (L5--L7) and the first
two lectures of Unit V (L8, regression foundations and least squares; L9,
maximum likelihood, $R^2$, model adequacy and over-fitting).
\textbf{Not examined} at mid-semester: the classification and margin losses of
Lecture~2 --- binary, categorical and sparse categorical cross-entropy, hinge
(SVM) loss and triplet loss --- and everything from Lecture~10 onwards
(logistic regression, multiple regression, regularisation). Adjusted $R^2$
with a given number of features $p$ \emph{is} examined; fitting a
multiple regression is not.

\textbf{What is inside.} Twenty-one topics, in lecture order. Each has
\textbf{the method} as a numbered recipe and \textbf{worked examples} done in
full; most also list the \textbf{mistakes} that cost marks and end with
\textbf{exercises} --- numerical ones in a purple box (answers at the back)
and short-answer ones with model answers in a grey box. The exercises draw on the course's shared
practice material (the chapter notes, the \emph{Practice \& Worked Solutions}
workbook and Class Test Set~B), recreated here with every answer checked. Two
complete practice papers with solutions close the booklet.

\textbf{How to use it.} Attempt every exercise on paper before you look at
Section~\ref{sec:answers}. Reading a solution feels like understanding; only
doing it is.
\end{keyidea}

\begin{center}
\small
\begin{tabular}{@{}rll@{\qquad}rll@{}}
\toprule
& \textbf{Topic} & \textbf{Lecture} & & \textbf{Topic} & \textbf{Lecture}\\ \midrule
1 & What ML is; types of learning & L1 & 10 & Outliers and transformations & L5\\
2 & Regression losses & L2 & 11 & Feature scaling & L6\\
3 & Gradient descent by hand & L3 & 12 & Encoding and feature engineering & L6\\
4 & Batch, stochastic, mini-batch & L3 & 13 & Data reduction and PCA & L6\\
5 & Momentum & L3 & 14 & Splitting and cross-validation & L7\\
6 & Adagrad and RMSprop & L4 & 15 & Imbalanced data & L7\\
7 & Adam & L4 & 16 & Data leakage and pipeline order & L5--L7\\
8 & Adadelta; comparing optimisers & L4 & 17 & Regression foundations & L8\\
9 & Missing data & L5 & 18 & Least squares by hand & L8\\
 & & & 19 & Direct regression; max.\ likelihood & L9\\
 & & & 20 & $R^2$, adjusted $R^2$, error summary & L9\\
 & & & 21 & Adequacy, over-fitting, CV & L9\\
\bottomrule
\end{tabular}
\end{center}

\begin{caution}[{{How numericals are marked}}]
A correct final number with no working earns little; a correct method with a
slip in one line earns most of the marks. Write the method visibly --- the
residual column, the gradient at each step, the accumulator, the sorted data
and the quartiles, the centred table --- carry \textbf{four decimal places},
and round only at the end. Finish every answer with a \textbf{check}:
residuals that sum to zero, a first Adam step of exactly $\eta$, a line that
passes through $(\bar x, \bar y)$.
\end{caution}

\begin{keyidea}[{{Conventions used throughout}}]
\begin{itemize}
  \item \textbf{Residual} $e = y - \hat y$ (actual minus predicted). Positive
    means the model \emph{under}-predicted.
  \item \textbf{Standard deviation.} Standardisation (\texttt{StandardScaler})
    uses the population $\sigma$, divisor $n$. The $z$-score outlier rule and
    the sample sd of a column use $s$, divisor $n - 1$.
  \item \textbf{Quartiles by hand}: sort, split into a lower and an upper
    half (for odd $n$ leave the median out of both), and take the median of
    each half (Tukey's hinges). NumPy's default \texttt{np.percentile}
    interpolates and can differ slightly; either is accepted if you say which
    you used.
  \item \textbf{Optimiser toy problem.} Unless stated otherwise the optimiser
    sections use two targets $y = 1$ and $y = 3$, a constant model $w$,
    half-squared row loss $\tfrac12(w - y)^2$ and their mean
    $J(w) = \tfrac12(w-2)^2 + \tfrac12$, so $J'(w) = w - 2$. With a loss
    $(w-c)^2$ (no half) the gradient is $2(w - c)$ --- say which you use.
  \item \textbf{Momentum} is $v \leftarrow \beta v + g$, $w \leftarrow w - \eta
    v$. Adam's first moment is $m \leftarrow \beta_1 m + (1-\beta_1) g$. Do not
    mix them.
\end{itemize}
\end{keyidea}

\tableofcontents
\newpage
\unit{Unit I --- Introduction to machine learning (L1)}
%======================================================================
\section{What machine learning is, and the three types of learning}
\label{sec:intro}

\method{naming and classifying a learning problem}
\begin{enumerate}
  \item \textbf{Definition} (Mitchell): a program learns from experience $E$
    on task $T$, measured by $P$, if its performance at $T$, as measured by
    $P$, improves with $E$. Name all three, and make $P$ specific.
  \item \textbf{Decide the type from the feedback.} Known answers in the
    training table $\Rightarrow$ \emph{supervised}; no answers, find structure
    $\Rightarrow$ \emph{unsupervised}; rewards for actions $\Rightarrow$
    \emph{reinforcement}.
  \item \textbf{If supervised, look at the target column}: a number on a
    continuous scale $\Rightarrow$ regression; a category $\Rightarrow$
    classification. The target decides, not the features.
  \item \textbf{Say what one training row contains}: the features, and the
    target if there is one.
  \item \textbf{Vocabulary}: a \emph{sample} is a row, a \emph{feature} a
    column used as input, the \emph{target} the column to predict. A
    \emph{parameter} is learned from data (a slope $b_1$); a
    \emph{hyperparameter} is set before learning (the $k$ of $k$-NN, a
    learning rate).
  \item \textbf{Judge generalisation on data the model never saw}: train on
    the training set, choose settings on the validation set, report once on
    the test set.
\end{enumerate}
\end{keyidea}

\begin{worked}[{{Example 1.1 --- classify the problem}}]
\small
\begin{tabular}{@{}p{5.2cm}p{3.3cm}p{5.6cm}@{}}
\toprule
\textbf{Problem} & \textbf{Type} & \textbf{One training row}\\ \midrule
Predict tomorrow's peak electricity demand in Dehradun & supervised,
regression & one day's weather and calendar features + that day's peak
demand\\[2pt]
Group 50{,}000 shoppers by purchasing behaviour, no categories given &
unsupervised (clustering) & one customer's purchase summary, no target\\[2pt]
Pneumonia or not, from 40{,}000 labelled X-rays & supervised,
classification & one X-ray's pixels + the radiologist's label\\[2pt]
Robot arm placing a cup: $+1$ success, $-1$ drop & reinforcement & no table
--- a stream of (state, action, reward)\\[2pt]
Predict a student's end-sem \emph{percentage} & supervised, regression &
attendance, internal marks, \dots\ + the percentage\\[2pt]
Predict the student's \emph{letter grade} A--F & supervised, classification
& the same features + the grade\\
\bottomrule
\end{tabular}

\normalsize
The last two rows are the same situation with the target recorded
differently --- which is exactly why the target column decides the type.

\textbf{$T$, $E$, $P$ for fraud detection.} $T$: label each card
transaction fraudulent or legitimate. $E$: past transactions with confirmed
outcomes. $P$: recall on the fraud class together with the false-alarm rate.
\emph{Not} plain accuracy: if $0.1\%$ of transactions are fraud, a model that
says ``legitimate'' every time is $99.9\%$ accurate and catches nothing.
\end{worked}

\begin{worked}[{{Example 1.2 --- 1-nearest neighbour, and reading results}}]
\textbf{1-NN by hand.} Five training flowers (petal length, petal width):
$(1.4, 0.2)$ setosa, $(1.3, 0.2)$ setosa, $(4.7, 1.4)$ versicolor,
$(4.5, 1.5)$ versicolor, $(6.0, 2.5)$ virginica. New flower $(4.6, 1.4)$.
Distances $\sqrt{\Delta_1^2 + \Delta_2^2}$: $3.42,\ 3.51,\ \mathbf{0.10},\
0.14,\ 1.78$. The nearest is $(4.7, 1.4)$, so 1-NN predicts
\textbf{versicolor}. It has no training step: it stores the data and does all
its work at prediction time, and it rests entirely on distance --- which is
why feature scaling (Section~\ref{sec:scaling}) matters to it.

\textbf{Reading train/test scores.} Three models on the same data:
\begin{center}
\small
\begin{tabular}{@{}lccc@{}}
\toprule
 & A & B & C\\ \midrule
training $R^2$ & $0.62$ & $0.88$ & $0.999$\\
test $R^2$ & $0.60$ & $0.85$ & $0.41$\\ \midrule
diagnosis & underfitting & good fit & overfitting\\
\bottomrule
\end{tabular}
\end{center}
A: both scores low and close --- too simple; add features or flexibility.
C: near-perfect on training, poor on test --- it has memorised. Deploy
\textbf{B}.

\textbf{Extrapolation.} A line fitted on houses of $1000$--$3000$ sq ft is
$\hat y = -50.9 + 0.0754x$ (lakh). At $2400$ sq ft: $130.1$ lakh. At $400$
sq ft: $-20.7$ lakh --- a negative price. Outside the range of the training
data the model has no information and gives no warning.
\end{worked}

\begin{caution}[{{Mistakes that cost marks}}]
\begin{itemize}
  \item \textbf{Deciding the type from the features} instead of the target.
  \item \textbf{A vague $P$} (``good performance''). Name the measure and say
    why it suits the task.
  \item \textbf{Quoting a training score as performance.} 100\% on the
    training data is consistent with memorisation. Ask for the test score.
  \item \textbf{Using ML where an exact rule exists} (GST on an invoice, a
    published fee formula): learning an approximation can only be worse.
\end{itemize}
\end{caution}

\begin{tryit}[{{Exercises 1 --- numerical}}]
\begin{enumerate}[label=\textbf{E1.\arabic*}]
  \item A model is $\hat y = 5 + 0.06x$ with $x$ in sq ft and $\hat y$ in lakh.
    Predict at $x = 1000$.
  \item An input matrix has 120 houses and 4 features. How many numerical
    feature entries does it hold? What is its shape, and the shape of $y$?
  \item Using the five flowers of Example 1.2, predict the species of
    $(1.5, 0.3)$ by 1-NN. Show all five distances.
  \item A hospital's readmission model is $94\%$ accurate overall, yet on the
    $6\%$ of patients who are readmitted it is right only $20\%$ of the time.
    How can both be true? What measure would have exposed it?
\end{enumerate}
\end{tryit}

\begin{concept}
\begin{enumerate}[label=\textbf{C1.\arabic*}]
  \item \emph{Traditional programming vs ML: what is input and what is
    output?} Traditional: data + hand-written rules in, answers out. ML: data
    + known answers in, the rules (a model) out.
  \item \emph{Classify: (a) group unlabelled houses by similarity; (b) a
    robot learning from rewards.} (a) Unsupervised --- no feedback at all,
    only structure. (b) Reinforcement --- feedback is a reward for an action,
    not a correct answer.
  \item \emph{Parameter vs hyperparameter, one example each.} A parameter is
    learned from data ($b_1$ in $\hat y = b_0 + b_1x$); a hyperparameter is set
    before learning and never changed by the algorithm ($k$ in $k$-NN, the
    learning rate $\eta$, a polynomial degree).
  \item \emph{Why is unsupervised learning harder to evaluate?} There is no
    correct answer to subtract the prediction from, so it is judged by
    stability, internal separation measures and usefulness --- all of which
    need judgement, not arithmetic.
  \item \emph{``fit'' vs ``predict''; which data for what?} \texttt{fit}
    learns parameters from the training set; \texttt{predict} applies them to
    new rows. Training set: fit. Validation set (or CV on training data):
    choose hyperparameters. Test set: one final, honest report.
  \item \emph{Removing a caste or religion column --- does the model then
    ignore it?} No. Locality, surname, school or pincode act as proxies, and a
    flexible model reconstructs the attribute from them; removing the column
    only removes your ability to audit the effect.
  \item \emph{Roles of NumPy, pandas and scikit-learn.} NumPy: fast
    vectorised arithmetic on arrays. pandas: loading, cleaning and grouping
    tables. scikit-learn: one \texttt{fit}/\texttt{predict} interface for
    preprocessing, models and evaluation.
  \item \emph{What is wrong with each line?}
    \texttt{model.fit(X\_test, y\_test)} --- fitting on the test set destroys
    it as a test; fit on the training data.
    \texttt{model.fit(df["Area"], df["Price"])} --- $X$ must be 2-D:
    \texttt{df[["Area"]]}.
    \texttt{X\_train, X\_test = train\_test\_split(X, y, ...)} --- with $X$
    and $y$ it returns four objects.
    \texttt{model.score(X\_train, y\_train)} --- a training score measures
    fit, not performance; score on the test set.
\end{enumerate}
\end{concept}

\unit{Unit II --- Loss functions (L2): the regression losses}
%======================================================================
\section{Regression losses: SSE, MSE, RMSE, MAE and Huber}
\label{sec:regloss}

\method{computing a regression loss}
\begin{enumerate}
  \item Write the residual row $e_i = y_i - \hat y_i$. Keep the sign: it
    tells you over- or under-prediction, and Huber needs $|e_i|$.
  \item One column per loss: $e_i^2$ (SSE, MSE, RMSE), $|e_i|$ (MAE), and for
    Huber first compare $|e_i|$ with $\delta$, \emph{then} pick the branch:
    \[
      L_\delta(e) = \begin{cases} \tfrac12 e^2 & |e| \le \delta\\
      \delta\bigl(|e| - \tfrac12\delta\bigr) & |e| > \delta\end{cases}
    \]
  \item $\mathrm{SSE} = \sum e_i^2$; $\mathrm{MSE} = \mathrm{SSE}/n$;
    $\mathrm{RMSE} = \sqrt{\mathrm{MSE}}$ (back in the units of $y$);
    $\mathrm{MAE} = \frac1n\sum|e_i|$; mean Huber $= \frac1n\sum L_\delta(e_i)$.
  \item Comment on which row dominates, and on robustness: squaring makes a
    doubled residual cost four times as much.
  \item \textbf{Gradients} (with respect to the prediction $\hat y$):
    $e^2 \to -2e$;\ $\tfrac12e^2 \to -e$;\ $|e| \to -\sgn(e)$ for $e \ne 0$;\
    Huber $\to -e$ inside, $-\delta\,\sgn(e)$ outside.
\end{enumerate}
\end{keyidea}

\begin{worked}[{{Example 2.1 --- three houses, every regression loss}}]
Observed $y = (5, 5, 5)$, predicted $\hat y = (4, 7, 2)$, Huber $\delta = 1$.
\begin{center}
\begin{tabular}{@{}lrrr|r@{}}
\toprule
house & 1 & 2 & 3 & $\Sigma$\\ \midrule
$e = y - \hat y$ & $1$ & $-2$ & $3$ & $2$\\
meaning & under & over & under & \\
$e^2$ & $1$ & $4$ & $9$ & $14$\\
$|e|$ & $1$ & $2$ & $3$ & $6$\\
$|e| \le 1$? & yes & no & no & \\
Huber & $\tfrac12(1)^2 = 0.5$ & $1(2 - 0.5) = 1.5$ & $1(3 - 0.5) = 2.5$ & $4.5$\\
\bottomrule
\end{tabular}
\end{center}
\[
  \mathrm{SSE} = 14, \quad \mathrm{MSE} = \tfrac{14}{3} = 4.6667, \quad
  \mathrm{RMSE} = 2.1602,
\]
\[
  \mathrm{MAE} = \tfrac63 = 2, \quad
  \text{Huber} = \tfrac{4.5}{3} = 1.5 .
\]
The mean \emph{signed} residual, $2/3$, is useless as an error measure:
residuals $(3, -3)$ would average $0$ while both miss by $3$. If $y$ is in
lakh, MSE is in lakh$^2$; RMSE and MAE are in lakh.

\textbf{Huber is smooth at the join.} At $e = \delta$ both branches give
$\tfrac12\delta^2$, and both slopes equal $\delta$. $\delta$ is measured in
the units of $y$ and sets where the linear tail begins.
\end{worked}

\begin{worked}[{{Example 2.2 --- the best constant: MSE picks the mean, MAE the median}}]
Targets $(0, 0, 9)$; the model must predict one constant $w$.
\begin{center}
\begin{tabular}{@{}lcc@{}}
\toprule
constant & SSE & sum of $|e|$\\ \midrule
$w = 3$ (mean) & $9 + 9 + 36 = \mathbf{54}$ & $3 + 3 + 6 = 12$\\
$w = 0$ (median) & $0 + 0 + 81 = 81$ & $0 + 0 + 9 = \mathbf{9}$\\
\bottomrule
\end{tabular}
\end{center}
Squared error is minimised by the mean ($3$); absolute error by the median
($0$). The single extreme target drags the MSE answer; the MAE answer stays
with the typical value.

\textbf{From loss to gradient.} Constant $w$, targets $(1, 3)$, row loss
$\tfrac12(w - y)^2$, at $w = 4$: row losses $4.5$ and $0.5$, so $J = 2.5$; row
gradients $w - y = 3$ and $1$, so $J'(4) = 2$. Loss and gradient are
different quantities --- the gradient is what the optimiser in Unit III uses.
Minimising SSE or MSE gives the same minimiser; their gradients differ by the
factor $n$, which matters for the learning rate.
\end{worked}

\begin{caution}[{{Mistakes that cost marks}}]
\begin{itemize}
  \item \textbf{Averaging signed residuals} as an error measure.
  \item \textbf{Huber's linear branch without $-\tfrac12\delta$}: it is
    $\delta(|e| - \tfrac12\delta)$, not $\delta|e|$.
  \item \textbf{Choosing the Huber branch from $e$ instead of $|e|$}: $e = -2$
    with $\delta = 1$ is \emph{outside}.
  \item \textbf{MSE in the wrong units.} Report RMSE when the units matter.
\end{itemize}
\end{caution}

\begin{tryit}[{{Exercises 2 --- numerical}}]
\begin{enumerate}[label=\textbf{E2.\arabic*}]
  \item $y = 50$, $\hat y = 47$: the residual? And for $y = 40$, $\hat y = 45$:
    is the model over- or under-predicting?
  \item Residuals $(-1, 2, 4)$: compute MSE, MAE and mean Huber loss with
    $\delta = 1$.
  \item With $\delta = 2$ and residuals $(1, -3)$, compute the mean Huber loss.
  \item Errors $(-4, 2)$: MAE and MSE. A residual changes from $3$ to $6$: by
    what factor does its squared penalty change? Its absolute penalty?
  \item $y = (10, 12, 15, 20)$, $\hat y = (11, 12, 13, 24)$. Compute MSE, MAE,
    RMSE and Huber ($\delta = 2$). Which row supplies most of the MSE?
  \item For $y = (1, 2, 2, 3, 12)$, find the best constant under MSE and under
    MAE. The $12$ is mistyped as $112$: what happens to each?
\end{enumerate}
\end{tryit}

\begin{concept}
\begin{enumerate}[label=\textbf{C2.\arabic*}]
  \item \emph{Is a loss function the same as an optimiser?} No. The loss
    states the penalty to be minimised; the optimiser is the rule that
    changes the parameters to reduce it.
  \item \emph{Show that RMSE $\ge$ MAE.} $\frac1n\sum e_i^2 -
    (\frac1n\sum|e_i|)^2$ is the variance of the numbers $|e_i|$, which is
    $\ge 0$; so $\mathrm{MSE} \ge \mathrm{MAE}^2$. Equality only when every
    $|e_i|$ is the same.
  \item \emph{When would you choose MAE or Huber over MSE?} When a few large
    residuals (outliers, or a heavy-tailed target) should not dominate the
    fit. MSE suits cases where a large miss really is disproportionately
    costly.
  \item \emph{Why does MAE have no ordinary derivative at $e = 0$?} $|e|$ has
    a corner there; any value in $[-1, 1]$ is a valid subgradient. Huber is
    differentiable everywhere, including at $\pm\delta$.
  \item \emph{Mean loss vs summed loss --- same training?} Same minimiser on a
    fixed dataset, but the gradient scales with $n$, so with a fixed learning
    rate (or an unscaled regularisation weight) the training behaves
    differently.
\end{enumerate}
\end{concept}
\unit{Unit III --- Optimizer functions (L3--L4)}
%======================================================================
\section{Gradient descent by hand}
\label{sec:gd}

\method{running gradient descent on paper}
\begin{enumerate}
  \item Write the objective and differentiate it. For the toy problem,
    row gradients are $g_i = w - y_i$ and the full gradient is their mean,
    $J'(w) = w - 2$. For the MSE of a line $\hat y = b + wx$,
    $\partial J/\partial w = \frac2n\sum(\hat y_i - y_i)x_i$ and
    $\partial J/\partial b = \frac2n\sum(\hat y_i - y_i)$.
  \item Update: $w_{\text{new}} = w_{\text{old}} - \eta\, J'(w_{\text{old}})$.
    Subtracting a positive gradient moves left; a negative one moves right.
  \item Tabulate $k$, $w_k$, $J(w_k)$, $J'(w_k)$ --- one row per step --- and
    check the loss falls.
  \item \textbf{Stability.} For the toy problem
    $w_{\text{new}} - 2 = (1 - \eta)(w_{\text{old}} - 2)$: the error shrinks
    when $|1 - \eta| < 1$, i.e.\ $0 < \eta < 2$. In general, for
    $J = \tfrac a2 (w - c)^2$ the factor is $1 - a\eta$ and the range is
    $0 < \eta < 2/a$; for $(w - c)^2$ (no half) it is $1 - 2\eta$ and
    $0 < \eta < 1$.
\end{enumerate}
\end{keyidea}

\begin{worked}[{{Example 3.1 --- the toy problem, and four learning rates}}]
Targets $1$ and $3$, start $w_0 = 6$, $\eta = 0.2$. At $w = 6$: row losses
$\tfrac12 5^2 = 12.5$ and $\tfrac12 3^2 = 4.5$, $J = 8.5$; row gradients $5$
and $3$, mean $4$.
\begin{center}
\begin{tabular}{@{}rrrr@{}}
\toprule
$k$ & $w_k$ & $J(w_k)$ & $J'(w_k) = w_k - 2$\\ \midrule
0 & 6.0000 & 8.5000 & 4.0000\\
1 & 5.2000 & 5.6200 & 3.2000\\
2 & 4.5600 & 3.7768 & 2.5600\\
3 & 4.0480 & 2.5972 & 2.0480\\
4 & 3.6384 & 1.8422 & \\
\bottomrule
\end{tabular}
\end{center}
The minimum is at $w = 2$ with $J(2) = 0.5$ --- not zero, because the two
targets disagree.

\begin{center}
\small
\begin{tabular}{@{}rrll@{}}
\toprule
$\eta$ & $1 - \eta$ & first steps from $6$ & behaviour\\ \midrule
$0.2$ & $0.8$ & $5.2,\ 4.56,\ 4.048$ & approaches gradually\\
$1$ & $0$ & $2$ & reaches the optimum in one step\\
$1.2$ & $-0.2$ & $1.2,\ 2.16,\ 1.968$ & overshoots, but converges\\
$2$ & $-1$ & $-2,\ 6,\ -2$ & undamped oscillation\\
\bottomrule
\end{tabular}
\end{center}
\end{worked}

\begin{worked}[{{Example 3.2 --- two parameters: a line through three points}}]
Fit $\hat y = b + wx$ to $x = (1,2,3)$, $y = (3,5,7)$ with
$J = \frac13\sum(\hat y_i - y_i)^2$, from $(b, w) = (0, 0)$, $\eta = 0.1$.

\textbf{Step 1.} Errors $\hat y - y = (-3, -5, -7)$, $J = 27.6667$.
\[
  \tfrac{\partial J}{\partial b} = \tfrac23(-15) = -10, \qquad
  \tfrac{\partial J}{\partial w} = \tfrac23\bigl(-3 - 10 - 21\bigr) = -22.6667,
\]
so $b = 1.0000$, $w = 2.2667$. \textbf{Step 2.} Errors
$(0.2667, 0.5333, 0.8000)$, $J = 0.3319$; gradient $(1.0667, 2.4889)$; so
$b = 0.8933$, $w = 2.0178$ and $J = 0.0053$. The true line is $y = 1 + 2x$.
Both partial derivatives are computed from the \emph{same} old parameters
before either is updated.
\end{worked}

\begin{caution}[{{Mistakes that cost marks}}]
\begin{itemize}
  \item \textbf{Adding the gradient} instead of subtracting it.
  \item \textbf{Confusing the loss value with the gradient.} At $w = 6$,
    $J = 8.5$ but $J' = 4$.
  \item \textbf{Dropping or doubling the factor 2.} With $\tfrac12(w-y)^2$ the
    row gradient is $w - y$; with $(w - y)^2$ it is $2(w - y)$.
  \item \textbf{``Larger $\eta$ is faster.''} $\eta = 0.8$ and $\eta = 1.2$
    have the same $|1 - \eta|$ and converge at the same rate.
\end{itemize}
\end{caution}

\begin{tryit}[{{Exercises 3 --- numerical}}]
\begin{enumerate}[label=\textbf{E3.\arabic*}]
  \item For $J(w) = \tfrac12(w - 3)^2$: find $J'(5)$, and the next parameter
    from $w = 5$ with $\eta = 0.1$.
  \item For that objective, derive the range of constant learning rates that
    converge to $3$ from any start, and the rate that gets there in one step.
  \item Toy problem (targets $1, 3$): at $w = 1$, find the full gradient and
    the direction of the step. From $w = 5$ with $\eta = 0.1$, take one step.
  \item Five steps of gradient descent on $f(w) = (w + 1)^2$ from $w_0 = 3$ with
    $\eta = 0.2$. By what factor does the distance to the minimum shrink?
  \item Fit $\hat y = b + wx$ to $x = (0,1,2)$, $y = (1,3,5)$ with
    $J = \frac13\sum(\hat y - y)^2$, from $(0,0)$, $\eta = 0.1$. Two steps;
    report $J$ after each.
\end{enumerate}
\end{tryit}

%======================================================================
\section{Batch, stochastic and mini-batch gradient descent}
\label{sec:sgd}

\method{updates, batches and epochs}
\begin{enumerate}
  \item An \textbf{update} is one change to the parameters; a \textbf{batch}
    is the set of rows that supplied its gradient; an \textbf{epoch} is one
    pass through the training data.
  \item \textbf{Full batch}: one update per epoch, the mean gradient over all
    rows. \textbf{SGD}: one row per update. \textbf{Mini-batch}: $B$ rows per
    update, averaged over the batch's actual size.
  \item Updates per epoch $= \lceil n/B\rceil$ when the last short batch is
    kept.
  \item In SGD each row's gradient is taken at the \emph{current} parameter,
    the one just produced --- not at the value the epoch started with.
\end{enumerate}
\end{keyidea}

\begin{worked}[{{Example 4.1 --- SGD and mini-batch on the toy problem}}]
Start $w = 6$, $\eta = 0.2$, rows visited in the order $1, 3, 1, 3$.
\begin{center}
\begin{tabular}{@{}ccrrr@{}}
\toprule
$t$ & selected $y$ & $w$ before & $g = w - y$ & $w$ after\\ \midrule
1 & 1 & 6.0000 & 5.0000 & 5.0000\\
2 & 3 & 5.0000 & 2.0000 & 4.6000\\
3 & 1 & 4.6000 & 3.6000 & 3.8800\\
4 & 3 & 3.8800 & 0.8800 & 3.7040\\
\bottomrule
\end{tabular}
\end{center}
\textbf{Mini-batch} with both rows: gradient $4 \Rightarrow w = 5.2$; then
$3.2 \Rightarrow w = 4.56$. (Here the batch is the whole toy dataset, so it is
also full-batch descent.)

\textbf{SGD can increase the full loss.} At the optimum $w = 2$ the full
gradient is $0$, but the row gradients are $+1$ and $-1$. One SGD step on the
$y = 1$ row moves $w$ to $2 - 0.2(1) = 1.8$ and the full loss rises from
$0.5$ to $0.52$. That is sampling noise, not a bug --- and why constant-rate
SGD fluctuates around the optimum.
\end{worked}

\begin{worked}[{{Example 4.2 --- one epoch of each, and counting updates}}]
$x = (1,2,3,4)$, $y = (3,5,8,9)$, model $\hat y = wx$, row loss
$(wx_i - y_i)^2$ with gradient $2x_i(wx_i - y_i)$, $w = 0$, $\eta = 0.02$.
Exact answer $w^\star = \sum xy/\sum x^2 = 73/30 = 2.4333$.
\begin{itemize}
  \item \textbf{Batch}: gradients at $0$ are $(-6, -20, -48, -72)$, mean
    $-36.5$; $w = 0.73$.
  \item \textbf{SGD}: $0.12,\ 0.5008,\ 1.2805,\ \mathbf{1.9010}$.
  \item \textbf{Mini-batch $B = 2$}: batch $\{1,2\}$ mean $-13$, $w = 0.26$;
    batch $\{3,4\}$ at $0.26$: $-43.32$ and $-63.68$, mean $-53.5$,
    $w = \mathbf{1.33}$.
\end{itemize}
Same work (four row-gradients), three answers: more updates per epoch mean
later gradients are taken at better parameters.

\textbf{Counting.} $100$ rows, $B = 20$: $5$ updates per epoch. $103$ rows,
$B = 20$, last batch kept: $\lceil 103/20\rceil = 6$ per epoch, $24$ in four
epochs. $n = 12{,}000$: $B = 1 \to 12{,}000$ updates, $B = 32 \to 375$,
$B = n \to 1$ --- and every one of them uses $12{,}000$ row-gradients per
epoch.
\end{worked}

\begin{caution}[{{Mistakes that cost marks}}]
\begin{itemize}
  \item \textbf{Evaluating every SGD gradient at the starting $w$.}
  \item \textbf{Summing a mini-batch gradient} instead of averaging it.
  \item \textbf{``Epoch'' $=$ ``update''.} An epoch contains $\lceil n/B\rceil$
    updates.
  \item \textbf{Shuffling time-series rows} as if they were independent.
\end{itemize}
\end{caution}

\begin{tryit}[{{Exercises 4 --- numerical}}]
\begin{enumerate}[label=\textbf{E4.\arabic*}]
  \item $100$ training rows, batch size $20$, no rows dropped: updates per
    epoch?
  \item Constant model $w = 4$, targets $(1, 3)$, row loss $\tfrac12(w - y)^2$:
    the batch-mean gradient?
  \item At $w = 2$ with targets $(1, 3)$, compare the full-data gradient with
    the single-row SGD gradient for $y = 1$. Can that SGD step increase the
    full loss?
  \item Toy problem from $w = 6$, $\eta = 0.2$: if the first selected target
    were $3$ rather than $1$, where would the first update land?
  \item $x = (2,1,3)$, $y = (4,3,5)$, $\hat y = wx$, $w_0 = 0$, $\eta = 0.05$.
    One batch step and one SGD epoch (in the order given). Which ends nearer
    $w^\star$?
  \item $n = 50{,}000$, $B = 250$, 10 epochs: how many updates, and how many
    row-gradient evaluations?
\end{enumerate}
\end{tryit}

%======================================================================
\section{Momentum}
\label{sec:momentum}

\method{momentum updates by hand}
\begin{enumerate}
  \item $v_t = \beta v_{t-1} + g_t$, then $w_t = w_{t-1} - \eta v_t$, from
    $v_0 = 0$.
  \item Tabulate $t$, $g_t$ (at the current $w$), $v_t$, $w_t$.
  \item Facts: $v_t = \sum_j \beta^j g_{t-j}$, a decaying sum of past
    gradients; with a constant gradient $g$, $v_t = g(1 - \beta^t)/(1-\beta)
    \to g/(1 - \beta)$, so the effective step is $\eta/(1-\beta)$ --- ten times
    $\eta$ at $\beta = 0.9$.
  \item The other convention, $v = \beta v + (1-\beta)g$, tends to $g$
    instead. Never mix them.
\end{enumerate}
\end{keyidea}

\begin{worked}[{{Example 5.1 --- history matters}}]
\textbf{One step.} $v_{\text{old}} = 2$, $g = -1$, $\beta = 0.9$, $\eta = 0.1$,
$w_{\text{old}} = 5$: $v_{\text{new}} = 0.9(2) + (-1) = 0.8$, and
$w_{\text{new}} = 5 - 0.1(0.8) = 4.92$. The current gradient says ``go
right'', but the stored velocity is still positive, so $w$ moves left.

\textbf{Constant gradient} $g = 2$, $\beta = 0.9$: $v = 2,\ 3.8,\ 5.42,\
6.878,\ 8.1902,\dots$ $= 2(1 - 0.9^t)/0.1 \to 20$.
\textbf{Alternating} $g = 2, -2, 2, -2$: $v = 2,\ -0.2,\ 1.82,\ -0.362$ ---
the contributions cancel. In a ravine the along-floor gradient is steady and
the across-wall gradient flips, so momentum speeds up along the floor and
damps the zig-zag.
\end{worked}

\begin{worked}[{{Example 5.2 --- four steps, two values of $\beta$}}]
$f(w) = (w - 2)^2$, $g = 2(w - 2)$, $w_0 = 6$, $\eta = 0.1$.
\begin{center}
\begin{tabular}{@{}rrrr@{}}
\toprule
$t$ & $g_t$ & $v_t = 0.5v_{t-1} + g_t$ & $w_t$\\ \midrule
1 & $8.0000$ & $8.0000$ & $5.2000$\\
2 & $6.4000$ & $10.4000$ & $4.1600$\\
3 & $4.3200$ & $9.5200$ & $3.2080$\\
4 & $2.4160$ & $7.1760$ & $\mathbf{2.4904}$\\
\bottomrule
\end{tabular}
\end{center}
Plain GD reaches $3.6384$ after four steps. With $\beta = 0.9$:
$v = 8,\ 13.6,\ 15.92,\ 14.824$ and $w = 5.2,\ 3.84,\ 2.248,\
\mathbf{0.7656}$ --- it has overshot the minimum, carried by nine tenths of
its past velocity although the gradient at step 4 is only $0.496$.
\end{worked}

\begin{caution}[{{Mistakes that cost marks}}]
\begin{itemize}
  \item \textbf{Using $g$ instead of $v$} in the parameter update.
  \item \textbf{Mixing conventions} ($\beta v + g$ versus $\beta v +
    (1-\beta)g$).
  \item \textbf{Calling an overshoot divergence.} Divergence means the swings
    grow.
\end{itemize}
\end{caution}

\begin{tryit}[{{Exercises 5 --- numerical}}]
\begin{enumerate}[label=\textbf{E5.\arabic*}]
  \item $v_0 = 0$, $\beta = 0.9$, $g_1 = g_2 = 2$: find $v_1$ and $v_2$.
  \item With $v_0 = 0$ and a constant gradient $g = 2$, derive $v_t$ for
    $\beta = 0.9$ and its limit.
  \item Three momentum steps on $f(w) = (w - 3)^2$ from $w_0 = 0$, $\eta = 0.1$,
    $\beta = 0.5$. Compare with plain GD ($1.4640$ after three steps).
  \item $\beta = 0.8$; gradients arrive as $4, 3, -1, -2$. Give the velocity
    after each.
\end{enumerate}
\end{tryit}

\begin{concept}
\begin{enumerate}[label=\textbf{C5.\arabic*}]
  \item \emph{What does $v$ store?} A decaying (unnormalised) sum of past
    gradients --- the accumulated direction of travel.
  \item \emph{What problem does momentum solve, and what does it cost?} It
    accelerates along directions where the gradient is consistent and damps
    oscillation where it flips; the price is overshoot, and one more
    hyperparameter $\beta$.
  \item \emph{Why does the convention matter for the learning rate?} With
    $v = \beta v + g$ the steady step is $\eta g/(1 - \beta)$; with $v = \beta
    v + (1-\beta)g$ it is $\eta g$. The same $\eta$ gives steps ten times
    apart at $\beta = 0.9$.
\end{enumerate}
\end{concept}
%======================================================================
\section{Adagrad and RMSprop}
\label{sec:adagrad}

\method{adaptive step sizes by hand}
\begin{enumerate}
  \item One accumulator per parameter, starting at $0$.
  \item \textbf{Adagrad} adds: $G_t = G_{t-1} + g_t^2$.
    \textbf{RMSprop} averages: $V_t = \rho V_{t-1} + (1-\rho) g_t^2$
    ($\rho = 0.9$ unless told otherwise).
  \item Step: $w_t = w_{t-1} - \eta\, g_t/(\sqrt{G_t} + \varepsilon)$ (or
    $\sqrt{V_t}$). Update the accumulator \emph{first}, then step. Ignore
    $\varepsilon$ by hand unless the accumulator is $0$.
  \item The numerator keeps the sign of $g$; only the denominator is squared.
  \item Checks: Adagrad's first step has size exactly $\eta$; RMSprop's is
    $\eta/\sqrt{1-\rho}$ ($3.1623\eta$ at $\rho = 0.9$).
\end{enumerate}
\end{keyidea}

\begin{worked}[{{Example 6.1 --- single steps, and three Adagrad steps}}]
\textbf{Adagrad, one step.} $G_{\text{old}} = 9$, $g = 4$, $\eta = 0.1$,
$w_{\text{old}} = 2$: $G_{\text{new}} = 9 + 16 = 25$, $\sqrt{G} = 5$, step
$= 0.1 \times 4/5 = 0.08$, $w_{\text{new}} = 1.92$.
\textbf{RMSprop, one step.} $V_{\text{old}} = 4$, $\rho = 0.9$, $g = 2$,
$\eta = 0.1$, $w_{\text{old}} = 5$: $V_{\text{new}} = 0.9(4) + 0.1(4) = 4$,
step $= 0.1 \times 2/2 = 0.1$, $w_{\text{new}} = 4.9$.

\textbf{Three Adagrad steps} on $f(w) = (w - 4)^2$, $g = 2(w - 4)$, $w_0 = 0$,
$\eta = 1$:
\begin{center}
\begin{tabular}{@{}rrrrrr@{}}
\toprule
$t$ & $g_t$ & $G_t$ & $\sqrt{G_t}$ & step & $w_t$\\ \midrule
1 & $-8.0000$ & $64.0000$ & $8.0000$ & $+1.0000$ & $1.0000$\\
2 & $-6.0000$ & $100.0000$ & $10.0000$ & $+0.6000$ & $1.6000$\\
3 & $-4.8000$ & $123.0400$ & $11.0923$ & $+0.4327$ & $2.0327$\\
\bottomrule
\end{tabular}
\end{center}
Step 1 is exactly $\eta$. Then the steps shrink faster than the gradient,
because $G$ only ever grows --- Adagrad's weakness on long runs. Its
strength: a rarely-active coordinate keeps a small $G$ and so a large step.
\end{worked}

\begin{worked}[{{Example 6.2 --- RMSprop on the same bowl}}]
$w_0 = 0$, $\eta = 0.5$, $\rho = 0.9$.
\begin{center}
\begin{tabular}{@{}rrrrrr@{}}
\toprule
$t$ & $g_t$ & $V_t = 0.9V_{t-1} + 0.1g_t^2$ & $\sqrt{V_t}$ & step & $w_t$\\ \midrule
1 & $-8.0000$ & $6.4000$ & $2.5298$ & $+1.5811$ & $1.5811$\\
2 & $-4.8377$ & $8.1004$ & $2.8461$ & $+0.8499$ & $2.4310$\\
3 & $-3.1380$ & $8.2750$ & $2.8766$ & $+0.5454$ & $2.9764$\\
\bottomrule
\end{tabular}
\end{center}
Check: $0.5/\sqrt{0.1} = 1.5811$. Adagrad at the same $\eta$ reaches only
$0.5,\ 0.8293,\ 1.0854$. RMSprop forgets old gradients (it remembers about
$1/(1-\rho) = 10$), so its denominator stays near the size of recent
gradients and can \emph{fall}. After $k$ zero gradients $V$ becomes
$\rho^k V_0$ --- the history changes although $w$ does not move.
\end{worked}

\begin{caution}[{{Mistakes that cost marks}}]
\begin{itemize}
  \item \textbf{Dividing by $G$ instead of $\sqrt{G}$.}
  \item \textbf{Forgetting $(1-\rho)$ in RMSprop}, or using $\rho g^2$.
  \item \textbf{Stepping with the old accumulator.} Update $G$ or $V$ first.
  \item \textbf{Losing the sign}: with $G_{\text{old}} = 0$, $g = -3$, $\eta =
    0.1$, the change is $-0.1(-3)/3 = +0.1$.
\end{itemize}
\end{caution}

\begin{tryit}[{{Exercises 6 --- numerical}}]
\begin{enumerate}[label=\textbf{E6.\arabic*}]
  \item $G_{\text{old}} = 9$, $g = 4$: the new accumulator? Then with
    $G_{\text{new}} = 25$, $g = 4$, $\eta = 0.1$, $w_{\text{old}} = 2$: $w_{\text{new}}$?
  \item $V_{\text{old}} = 4$, $\rho = 0.9$, $g = 2$, $\eta = 0.1$,
    $w_{\text{old}} = 5$: $w_{\text{new}}$ by RMSprop. And for $V_{\text{old}} =
    0$, $g = 2$: $V_{\text{new}}$ and the denominator.
  \item Three Adagrad steps on $f(w) = (w - 1)^2$ from $w_0 = 5$, $\eta = 0.5$.
  \item Gradients $4, 2, 1, 0.5$: the Adagrad and the RMSprop ($\rho = 0.9$)
    accumulators after each, and $g/\sqrt{\cdot}$ for each. Which ratio
    stays larger, and why?
  \item After $k$ consecutive zero gradients, what is $V$ if it started at
    $V_0$? Does the parameter move during those updates?
\end{enumerate}
\end{tryit}

%======================================================================
\section{Adam, with bias correction}
\label{sec:adam}

\method{Adam by hand}
\begin{enumerate}
  \item Defaults $\beta_1 = 0.9$, $\beta_2 = 0.999$; $m_0 = v_0 = 0$; $t$
    counts updates from $1$ and is never reset at an epoch boundary.
  \item $m_t = \beta_1 m_{t-1} + (1-\beta_1) g_t$ (signed history);
    $v_t = \beta_2 v_{t-1} + (1-\beta_2) g_t^2$ (squared history).
  \item Bias correction: $\hat m_t = m_t/(1 - \beta_1^t)$,
    $\hat v_t = v_t/(1 - \beta_2^t)$. Divisors: $0.1$ and $0.001$ at $t = 1$;
    $0.19$ and $0.001999$ at $t = 2$.
  \item Step: $w_t = w_{t-1} - \eta\,\hat m_t/(\sqrt{\hat v_t} + \varepsilon)$.
    Carry the \emph{uncorrected} $m_t$, $v_t$ forward.
  \item Check: $\hat m_1 = g_1$ and $\hat v_1 = g_1^2$, so the first step is
    exactly $\pm\eta$, whatever the size of $g_1$.
\end{enumerate}
\end{keyidea}

\begin{worked}[{{Example 7.1 --- Adam, two steps from a given gradient sequence}}]
$w_0 = 5$, $\eta = 0.1$, gradients $g_1 = 2$, $g_2 = 1$.

\textbf{Step 1.} $m_1 = 0.1(2) = 0.2$; $v_1 = 0.001(4) = 0.004$;
$\hat m_1 = 0.2/0.1 = 2$; $\hat v_1 = 0.004/0.001 = 4$; direction
$2/\sqrt4 = 1$; $w_1 = 5 - 0.1 = 4.9$.

\textbf{Step 2} (keep the history). $m_2 = 0.9(0.2) + 0.1(1) = 0.28$;
$v_2 = 0.999(0.004) + 0.001(1) = 0.004996$;
$\hat m_2 = 0.28/0.19 = 1.473684$; $\hat v_2 = 0.004996/0.001999 = 2.499250$;
direction $1.473684/1.580902 = 0.932180$; $w_2 = 4.9 - 0.0932 = 4.806782$.

The second direction is not simply $\sgn(g_2) = 1$: both stored histories
still remember $g_1 = 2$.
\end{worked}

\begin{worked}[{{Example 7.2 --- three steps on a bowl, and without the correction}}]
$f(w) = (w - 4)^2$, $w_0 = 0$, $\eta = 0.5$.
\begin{center}
\begin{tabular}{@{}rrrrrrr@{}}
\toprule
$t$ & $g_t$ & $m_t$ & $v_t$ & $\hat m_t$ & $\sqrt{\hat v_t}$ & $w_t$\\ \midrule
1 & $-8.0000$ & $-0.8000$ & $0.064000$ & $-8.0000$ & $8.0000$ & $0.5000$\\
2 & $-7.0000$ & $-1.4200$ & $0.112936$ & $-7.4737$ & $7.5164$ & $0.9972$\\
3 & $-6.0057$ & $-1.8786$ & $0.148891$ & $-6.9320$ & $7.0484$ & $1.4889$\\
\bottomrule
\end{tabular}
\end{center}
Steps $0.500,\ 0.497,\ 0.492$: Adam moves at almost exactly the speed $\eta$.

\textbf{Without bias correction} the first step would be
$0.5 \times 0.8/\sqrt{0.064} = 1.5811$ --- three times $\eta$ --- because $m$ is
shrunk by $1 - \beta_1 = 0.1$ but $\sqrt v$ only by $\sqrt{1-\beta_2} =
0.0316$; the ratio is inflated by $0.1/0.0316 = 3.1623$.
\end{worked}

\begin{caution}[{{Mistakes that cost marks}}]
\begin{itemize}
  \item \textbf{$1 - \beta$ instead of $1 - \beta^t$} at step 2 or later.
  \item \textbf{Carrying $\hat m$, $\hat v$ forward} instead of $m$, $v$.
  \item \textbf{$v = m^2$.} $v$ averages $g^2$.
  \item \textbf{Calling $v$ the variance.} It is an uncentred second moment.
\end{itemize}
\end{caution}

\begin{tryit}[{{Exercises 7 --- numerical}}]
\begin{enumerate}[label=\textbf{E7.\arabic*}]
  \item $t = 1$, $g = 2$, $w_0 = 5$, $\eta = 0.1$, defaults: find $w_1$. If
    instead $g_1 = -2$: $\hat m_1$, $\hat v_1$ and the signed step?
  \item Two Adam steps on $f(w) = (w - 1)^2$ from $w_0 = 3$, $\eta = 0.1$.
  \item Show that the corrected first step is $-\eta\,\sgn(g_1)$. What is it
    for $g_1 = -0.003$, $\eta = 0.01$?
  \item Without correction, how large is the first step (as a multiple of
    $\eta$) for $\beta_2 = 0.999$? For $\beta_2 = 0.99$ (with $\beta_1 = 0.9$)?
\end{enumerate}
\end{tryit}

%======================================================================
\section{Adadelta, and comparing optimisers}
\label{sec:adadelta}

\method{Adadelta by hand}
\begin{enumerate}
  \item Two decaying averages: $G$ of squared gradients and $U$ of squared
    \emph{updates}, both from $0$; $\rho = 0.9$; $\varepsilon$ inside both
    roots.
  \item In this order: (i) $G_t = \rho G_{t-1} + (1-\rho)g_t^2$;
    (ii) $\Delta w_t = -\dfrac{\sqrt{U_{t-1} + \varepsilon}}{\sqrt{G_t +
    \varepsilon}}\, g_t$; (iii) $w_t = w_{t-1} + \Delta w_t$;
    (iv) $U_t = \rho U_{t-1} + (1-\rho)\Delta w_t^2$.
  \item The numerator uses the \emph{previous} $U_{t-1}$ --- $U_t$ cannot
    exist until $\Delta w_t$ does. There is no learning rate in the original
    form: the ratio has the units of $w$.
\end{enumerate}
\end{keyidea}

\begin{worked}[{{Example 8.1 --- one step with history, and the first step from zero}}]
\textbf{With history.} $G_{\text{old}} = 4$, $U_{\text{old}} = 1$, $\rho =
0.9$, $g = 2$, $\varepsilon$ ignored: $G = 0.9(4) + 0.1(4) = 4$; ratio
$\sqrt1/\sqrt4 = 0.5$; $\Delta w = -0.5 \times 2 = -1$; $U = 0.9(1) + 0.1(1)
= 1$.

\textbf{From zero} on $f(w) = (w - 3)^2$, $w_0 = 0$, $\varepsilon = 10^{-6}$:
$g_1 = -6$, $G_1 = 0.1 \times 36 = 3.6$, $\sqrt{G_1 + \varepsilon} = 1.8974$;
$\sqrt{U_0 + \varepsilon} = \sqrt{10^{-6}} = 0.001$; so
$\Delta w_1 = -(0.001/1.8974)(-6) = +0.003162$. Tiny: Adagrad's first step
here would be $\eta$, Adadelta's is set by $\varepsilon$. Then $U_1 = 0.1
(0.003162)^2 = 10^{-6}$, the numerator grows to $\sqrt{2\times10^{-6}} =
0.001414$, and the steps grow as updates accumulate.
\end{worked}

\begin{worked}[{{Example 8.2 --- what each optimiser stores}}]
\small
\begin{tabular}{@{}lll@{}}
\toprule
\textbf{Method} & \textbf{State beyond $w$} & \textbf{Main sensitivity}\\ \midrule
Plain SGD / batch & none & learning rate, sampling, feature scaling\\
Momentum & decaying signed sum $v$ & $\beta$ and the velocity convention\\
Adagrad & cumulative $g^2$ per coordinate & ever-growing denominator\\
RMSprop & decaying mean of $g^2$ & $\rho$, placement of $\varepsilon$\\
Adam & $m$, $v$ and the update count $t$ & bias correction, variant\\
Adadelta & decaying means of $g^2$ and $\Delta w^2$ & update order, $\varepsilon$\\
\bottomrule
\end{tabular}

\normalsize
\textbf{A fair comparison} keeps the model, loss reduction, data split,
preprocessing and metric fixed; gives each optimiser its own tuned learning
rate (the same number means different step sizes in different rules);
repeats stochastic runs; chooses on validation data; and reports on untouched
test data together with the compute used. A lower training loss is not proof
of better generalisation.
\end{worked}

\begin{tryit}[{{Exercises 8 --- numerical}}]
\begin{enumerate}[label=\textbf{E8.\arabic*}]
  \item $G_{\text{old}} = 4$, $U_{\text{old}} = 1$, $\rho = 0.9$, $g = 2$,
    $\varepsilon$ ignored: compute $\Delta w$.
  \item $G_{\text{old}} = 1$, $U_{\text{old}} = 0.25$, $\rho = 0.9$, $g = -2$,
    $\varepsilon$ ignored: compute $G$, $\Delta w$ and the new $U$.
\end{enumerate}
\end{tryit}

\begin{concept}
\begin{enumerate}[label=\textbf{C8.\arabic*}]
  \item \emph{What do Adam's $m$ and $v$ average?} $m$: the signed gradients
    (direction, as in momentum). $v$: the squared gradients (scale, as in
    RMSprop).
  \item \emph{Why can Adagrad learn rare coordinates differently, and why can
    it stall?} Each coordinate has its own $G$; rarely active ones keep a
    small $G$ and a large step. But $G$ never decreases, so every step shrinks
    towards zero on long runs.
  \item \emph{What distinguishes RMSprop from Adagrad?} An exponentially
    weighted \emph{average} of $g^2$ instead of a cumulative sum, so old
    gradients lose weight and the denominator can fall.
  \item \emph{Which two moving averages does Adadelta track?} $E[g^2]$ and
    $E[\Delta w^2]$.
  \item \emph{Two correct demonstrations give different trajectories --- why?}
    Different loss reduction (sum or mean, with or without $\tfrac12$),
    momentum convention, $\varepsilon$ placement, initial values or learning
    rate. Compare conventions before declaring an error.
  \item \emph{Your analytic gradient is twice the finite-difference value.
    What do you check first?} Whether the loss has a $\tfrac12$ (or an
    averaging factor) that the derivative dropped.
\end{enumerate}
\end{concept}
\unit{Unit IV --- Data preprocessing (L5--L7)}

\begin{keyidea}[{{The housing dataset used in Unit IV}}]
Twelve houses offered for sale in Dehradun (from the course's chapter notes).
Area in sq ft, Age in years, Dist in km to the city centre, buyer Income and
Price in lakh rupees.
\begin{center}
\scriptsize
\setlength{\tabcolsep}{3.5pt}
\begin{tabular}{@{}lrrrrllllrrr@{}}
\toprule
ID & Area & Bed & Bath & Age & Location & Garage & Condition & & Dist & Income & Price\\ \midrule
H01 & 1000 & 2 & 1 & 12 & Clement Town & No & Average & & 8.5 & 9.5 & 42\\
H02 & 1500 & 3 & 2 & 5 & Rajpur Road & Yes & Good & & 3.2 & 18 & 78\\
H03 & ? & 3 & 2 & 9 & Sahastradhara & No & Average & & 8.0 & 14 & 61\\
H04 & 2000 & 4 & 3 & 2 & Rajpur Road & Yes & Excellent & & 2.5 & 32 & 120\\
H05 & 1200 & 2 & 1 & 18 & Prem Nagar & No & Poor & & 12.4 & 7 & 38\\
H06 & 2500 & 4 & 3 & 7 & Rajpur Road & Yes & Good & & 4.0 & 28 & 140\\
H07 & 1800 & 3 & 2 & ? & Clement Town & Yes & Good & & 7.8 & 16 & 72\\
H08 & 3000 & 5 & 4 & 3 & Sahastradhara & Yes & Excellent & & 5.5 & 45 & 185\\
H09 & 1350 & 3 & 1 & 15 & Prem Nagar & ? & Average & & 13.0 & 11 & 46\\
H10 & 2200 & 4 & 3 & 6 & Clement Town & Yes & Good & & 6.5 & 24 & 105\\
H11 & 9500 & 6 & 6 & 4 & Rajpur Road & Yes & Excellent & & 2.0 & 95 & 620\\
H12 & 1600 & 3 & 2 & 150 & Prem Nagar & No & Average & & 10.2 & ? & 58\\
\bottomrule
\end{tabular}
\end{center}
Seven problems hide in it: three blanks (H03 Area, H07 Age, H09 Garage) and a
fourth (H12 Income); an impossible value (H12 Age $= 150$); a genuine extreme
(H11); columns on wildly different scales; three text columns; an identifier
(ID) that carries no information; and Income, which is recorded \emph{after}
the sale --- a leak (Section~\ref{sec:leakage}).
\end{keyidea}

%======================================================================
\section{Missing data}
\label{sec:missing}

\method{handling a blank}
\begin{enumerate}
  \item \textbf{Ask why it is missing.} MCAR: unrelated to anything. MAR:
    depends on other \emph{observed} columns. MNAR: depends on the missing
    value itself --- cannot be detected from the data, only from how it was
    collected.
  \item \textbf{Fix known errors first}, then compute any statistic you will
    fill with.
  \item \textbf{Choose}: delete rows (tiny fraction, MCAR only); delete a
    column (mostly empty); impute the \emph{median} for a skewed numeric column,
    the \emph{mean} for a symmetric one, the \emph{mode} for a categorical one
    --- group-wise when a sensible group exists; add a
    \emph{missing-indicator} column when missingness may itself be
    informative (MNAR); or a model-based imputer (KNN, iterative).
  \item \textbf{Compute from observed values only}, on the \textbf{training}
    rows only.
  \item Know the side effect: mean imputation leaves the mean unchanged but
    shrinks the spread --- the sample variance is multiplied by
    $(m-1)/(n-1)$ when $m$ of $n$ values are observed.
\end{enumerate}
\end{keyidea}

\begin{worked}[{{Example 9.1 --- the blanks in the housing table}}]
\textbf{Deletion.} Blanks in H03, H07, H09 and H12: $12 \to 8$ rows, $66.7\%$
kept --- a third of the data thrown away to repair four cells of $132$.

\textbf{H03 Area (numeric).} Eleven observed values, sum $27{,}650$:
mean $= 27650/11 = 2513.64$. Sorted: $1000, 1200, 1350, 1500, 1600,
\mathbf{1800}, 2000, 2200, 2500, 3000, 9500$; median $=$ 6th value $= 1800$.
The single 9500 pulls the mean $40\%$ above the median; nine of eleven houses
are below the mean. Fill with the \textbf{median, 1800} (also plausible for a
3-bed, 2-bath house).

\textbf{H09 Garage (categorical).} Global mode: Yes (7) vs No (4) $\Rightarrow$
Yes. But the other Prem Nagar houses (H05, H12) have no garage: the
\textbf{group-wise mode} within Location says No --- more likely correct.

\textbf{H07 Age, and the order of operations.} With H12's $150$ still in: mean
$231/11 = 21.00$, median $7$. After correcting $150 \to 15$: mean $96/11 =
8.73$, median $7$. Impute before fixing the error and H07 is recorded as 21
years old; the other Clement Town houses are 12 and 6. Fix errors first --- and note that the
median did not move.

\textbf{H12 Income, with an indicator.} Fill with the median ($18$) and add
\texttt{Income\_missing} $= 1$: a column is added, no row is. (Income is later
removed as a leak --- Section~\ref{sec:leakage}.)
\end{worked}

\begin{worked}[{{Example 9.2 --- mean imputation shrinks the spread}}]
Marks $40, 46, ?, 50, 54, ?, 60, 98$. Observed ($m = 6$): mean $348/6 = 58$,
median $(50 + 54)/2 = 52$, sample sd $20.7461$. Right-skewed (mean $>$
median) $\Rightarrow$ the median is the safer fill.

If both holes are filled with the mean $58$: the mean stays $58$, the sum of
squared deviations is unchanged, but $n$ grows from $6$ to $8$, so
$s^2_{\text{after}} = \tfrac57 s^2_{\text{before}}$ and
$s_{\text{after}} = 20.7461 \times 0.8452 = 17.5336$. The model sees a spread
$15\%$ narrower than the real one, and the column's relationship with the
target is weakened.
\end{worked}

\begin{caution}[{{Mistakes that cost marks}}]
\begin{itemize}
  \item \textbf{Counting the blank} in $n$.
  \item \textbf{Imputing before correcting known errors.}
  \item \textbf{Computing the fill value on all rows} before splitting ---
    leakage.
  \item \textbf{Filling a blank that is not missing}: an empty
    \texttt{Garage\_Area} may simply mean $0$.
\end{itemize}
\end{caution}

\begin{tryit}[{{Exercises 9 --- numerical}}]
\begin{enumerate}[label=\textbf{E9.\arabic*}]
  \item Training values $10, 20, ?, 40$. Impute with the observed median.
  \item A column holds $1200, 1500, ?, 1800, 2000, 2500, 9000$. Mean and median
    of the observed values; which do you impute with?
  \item The eleven observed buyer incomes are $7, 9.5, 11, 14, 16, 18, 24, 28,
    32, 45, 95$. Mean and median; which fills H12's blank, and why?
  \item $n = 100$ with $20$ values mean-imputed. By what factor does the
    population sd shrink? The sample sd?
\end{enumerate}
\end{tryit}

\begin{concept}
\begin{enumerate}[label=\textbf{C9.\arabic*}]
  \item \emph{Why is the median usually safer than the mean? When is the mean
    as good?} The mean uses magnitudes, so one extreme value moves it without
    limit; the median uses ranks. For a roughly symmetric column the two
    nearly coincide.
  \item \emph{Distinguish MCAR, MAR and MNAR. Which cannot be detected from
    the data?} See the method box. MNAR cannot, because the evidence needed is
    exactly what is missing; handle it with an indicator column and state the
    limitation.
  \item \emph{Missing income is more common among high earners. Deletion or
    median imputation?} This is MNAR. Deleting drops the high earners and
    biases the data; plain median imputation understates their income. Impute
    the median \emph{and} add \texttt{Income\_missing}, fitted on training
    data only, and report the limitation.
  \item \emph{Is every unusual house price a data error?} No --- see
    Section~\ref{sec:outliers}.
\end{enumerate}
\end{concept}

%======================================================================
\section{Outliers and transformations}
\label{sec:outliers}

\method{detecting and treating an extreme value}
\begin{enumerate}
  \item \textbf{Sort}, find $Q_1$ and $Q_3$ (hinges: medians of the lower and
    upper halves), $\IQR = Q_3 - Q_1$.
  \item \textbf{Tukey fences}: flag values below $Q_1 - 1.5\IQR$ or above
    $Q_3 + 1.5\IQR$.
  \item \textbf{$z$-score}: $z = (x - \bar x)/s$ (sample $s$), flag
    $|z| > 3$. Beware \emph{masking} --- the outlier inflates the $\bar x$ and
    $s$ it is judged by --- and the ceiling $|z| \le (n - 1)/\sqrt n$.
  \item \textbf{Modified $z$}: $M = 0.6745(x - \tilde x)/\MAD$ with $\MAD =
    \operatorname{median}|x_i - \tilde x|$, flag $|M| > 3.5$.
  \item \textbf{Decide}, don't delete by reflex. Impossible $\Rightarrow$
    correct it or treat as missing. Genuine $\Rightarrow$ keep it and use a
    tolerant method (log transform, robust scaler, MAE/Huber, trees), clip
    (winsorise) at the fence, or scope the model and say so.
  \item \textbf{Shape}: skewness $g_1 = m_3/m_2^{3/2}$ with $m_k =
    \frac1n\sum(x_i - \bar x)^k$; $g_1 > 0$ means a long right tail --- a
    candidate for $\ln x$.
\end{enumerate}
\end{keyidea}

\begin{worked}[{{Example 10.1 --- the twelve prices}}]
Sorted: $38, 42, 46, 58, 61, 72 \mid 78, 105, 120, 140, 185, 620$.

\textbf{IQR.} $Q_1 = (46 + 58)/2 = 52$, $Q_3 = (120 + 140)/2 = 130$,
$\IQR = 78$. Fences $52 - 117 = -65$ and $130 + 117 = 247$. $620 > 247$:
\textbf{H11 flagged}. H08 (185) is not --- the rule picks out a separate
observation, not merely the top of the range. (NumPy's interpolated quartiles
are $55$ and $125$, upper fence $230$: the same verdict.)

\textbf{$z$-score.} $\bar x = 130.42$, $s = 160.46$:
$z = (620 - 130.42)/160.46 = 3.05$ --- only just over $3$, because H11 itself
inflated $s$ from $46.63$ to $160.46$. Judged against the other eleven houses
($\bar x = 85.91$, $s = 46.63$) it scores $z = 11.45$. And with $n = 12$ no
point can exceed $11/\sqrt{12} = 3.175$: the $|z| > 3$ rule could never fire
on ten points or fewer.

\textbf{Decision.} H11 is a genuine luxury property: keep it if you will price
such houses, or restrict the model's stated domain and say so. H12's Age $=
150$ is impossible for its context: correct it to $15$.
\textbf{Treatments.} Clipping sets $620 \to 247$ (a fiction, but capped).
$\ln$: prices $38$--$620$ (ratio $16.3\times$) become $3.64$--$6.43$
(ratio $1.77\times$); nothing is deleted, the scale is changed, and
predictions must be exponentiated back.
\end{worked}

\begin{worked}[{{Example 10.2 --- skewness, and why the log helps}}]
$x = (1, 1, 2, 4, 7)$, $\bar x = 3$, deviations $(-2,-2,-1,1,4)$:
$m_2 = 26/5 = 5.2$, $m_3 = (-8 - 8 - 1 + 1 + 64)/5 = 9.6$,
$g_1 = 9.6/5.2^{1.5} = 9.6/11.8579 = +0.8096$ --- a right tail, produced by
the single $7$.

$\ln$ of the geometric sequence $1, 4, 16, 64$ gives $0, 1.3863, 2.7726,
4.1589$: equal gaps of $\ln 4$. The log turns ``four times bigger'' into
``one step bigger'', which is why prices, areas and incomes usually want it.
\end{worked}

\begin{caution}[{{Mistakes that cost marks}}]
\begin{itemize}
  \item \textbf{``Outlier $=$ error.''} H11 is rare but true.
  \item \textbf{Concluding ``no outliers'' from $|z| \le 3$} on a small sample.
  \item \textbf{Recomputing the fences on the test set.} $Q_1$, $Q_3$ and the
    fences are learned from training data and applied unchanged.
  \item \textbf{Deleting rows from the test set.} It must reflect reality.
\end{itemize}
\end{caution}

\begin{tryit}[{{Exercises 10 --- numerical}}]
\begin{enumerate}[label=\textbf{E10.\arabic*}]
  \item $Q_1 = 10$, $Q_3 = 18$: the IQR and the upper fence.
  \item Apply the IQR rule to $20, 25, 30, 35, 40, 45, 50, 200$: $Q_1$, $Q_3$,
    both fences, flagged values.
  \item $12, 14, 15, 15, 16, 17, 18, 19, 60$ ($n = 9$). Apply the IQR rule
    (leave the median out of both halves), the $z$-score rule and the
    modified $z$ to the $60$. Why does the $z$ rule fail?
  \item The largest possible $|z|$ for $n = 10$? $n = 11$? For which can
    $|z| > 3$ ever fire?
  \item Skewness $g_1$ of $(0, 0, 0, 0, 10)$. Compare with the ceiling
    $(n - 2)/\sqrt{n - 1}$.
\end{enumerate}
\end{tryit}

\begin{concept}
\begin{enumerate}[label=\textbf{C10.\arabic*}]
  \item \emph{A student removed every row with $|z| > 3$, then mean-imputed,
    then split. Three mistakes?} (i) Deleting by a rule instead of deciding
    error vs genuine (and the $z$ rule is masked on small samples);
    (ii) mean imputation of skewed columns, computed on all rows; (iii) every
    statistic was learned before the split --- leakage into the test set.
  \item \emph{What changes when prices are log-transformed before
    regression?} The model predicts $\ln(\text{price})$; errors become
    roughly proportional (percentage) errors; predictions must be
    exponentiated back, and an RMSE on the log scale is not in lakh.
  \item \emph{Where does the 1.5 come from?} A convenient compromise: for
    normal data the fences cover about $99.3\%$, so about $0.7\%$ of clean
    points are flagged. Use $3$ to flag only very extreme points.
\end{enumerate}
\end{concept}

%======================================================================
\section{Feature scaling}
\label{sec:scaling}

\method{the four scalers}
\begin{enumerate}
  \item From the \textbf{training} column compute the statistics, then apply
    $x \mapsto (x - a)/b$ to training \emph{and} test rows:
    \begin{center}\small
    \begin{tabular}{@{}llll@{}}
    \toprule
    scaler & $a$ & $b$ & result\\ \midrule
    min--max (normalisation) & $x_{\min}$ & $x_{\max} - x_{\min}$ & $[0, 1]$ on training data\\
    standardisation ($z$) & $\bar x$ & $\sigma$ (divisor $n$) & mean $0$, sd $1$\\
    robust & median & $\IQR$ & median $0$, IQR $1$\\
    max-abs & $0$ & $\max|x|$ & $[-1, 1]$, zeros kept\\
    \bottomrule
    \end{tabular}
    \end{center}
  \item Test values may land outside $[0, 1]$: that is correct, and it
    signals a value outside the training range. Never refit on the test set.
  \item \textbf{Who needs it?} Anything that compares or combines features
    numerically: $k$-NN, $k$-means, SVM (RBF), neural networks, gradient
    descent, ridge/lasso. Not trees or forests (they split one feature at a
    time on order), and not least squares by the exact formula.
\end{enumerate}
\end{keyidea}

\begin{worked}[{{Example 11.1 --- scaling changes the nearest neighbour}}]
Query house (2000 sq ft, 4 bed); A $=$ (2050, 2); B $=$ (1500, 4).

\textbf{Unscaled}: $d(Q, A) = \sqrt{50^2 + 2^2} = 50.04$;
$d(Q, B) = \sqrt{500^2 + 0} = 500$. A is ``nearest'' --- the bedroom
difference supplies $4$ of $2504$, $0.16\%$ of the squared distance.

\textbf{Min--max} with Area $\in [1000, 3000]$, Bedrooms $\in [2, 5]$:
Q $= (0.500, 0.667)$, A $= (0.525, 0)$, B $= (0.250, 0.667)$;
$d(Q, A) = \sqrt{0.000625 + 0.4444} = 0.667$, $d(Q, B) = 0.250$. Now
\textbf{B is nearest}: a 4-bedroom house is priced from a 4-bedroom house.
Scaling added no information; it stopped one column monopolising the
distance.

\textbf{Min--max is ruined by an outlier.} With H11 in the column
($x_{\max} = 9500$), a 2000 sq ft house scales to $1000/8500 = 0.1176$ and a
3000 sq ft house to $0.2353$ --- every ordinary house is squashed into the
bottom quarter.
\end{worked}

\begin{worked}[{{Example 11.2 --- one column through four scalers}}]
$v = (2, 4, 4, 4, 5, 5, 7, 9)$: mean $5$, deviations $(-3,-1,-1,-1,0,0,2,4)$,
$\sigma = \sqrt{32/8} = 2$; min $2$, max $9$; median $4.5$; hinges
$Q_1 = 4$ (median of $2,4,4,4$), $Q_3 = 6$ (median of $5,5,7,9$), $\IQR = 2$.
\begin{center}
\small
\begin{tabular}{@{}lrrrrrrrr@{}}
\toprule
$v$ & 2 & 4 & 4 & 4 & 5 & 5 & 7 & 9\\ \midrule
standard $(v - 5)/2$ & $-1.5$ & $-0.5$ & $-0.5$ & $-0.5$ & $0$ & $0$ & $1$ & $2$\\
min--max $(v - 2)/7$ & $0$ & $.2857$ & $.2857$ & $.2857$ & $.4286$ & $.4286$ & $.7143$ & $1$\\
robust $(v - 4.5)/2$ & $-1.25$ & $-0.25$ & $-0.25$ & $-0.25$ & $0.25$ & $0.25$ & $1.25$ & $2.25$\\
max-abs $v/9$ & $.2222$ & $.4444$ & $.4444$ & $.4444$ & $.5556$ & $.5556$ & $.7778$ & $1$\\
\bottomrule
\end{tabular}
\end{center}
(\texttt{RobustScaler} uses NumPy quartiles, $4$ and $5.5$, so it divides by
$1.5$; by hand, state your rule.) \textbf{A test value outside the range:}
training Area $200$--$800$, test house $950$: correctly $(950 - 200)/600 =
1.25$. A scaler wrongly fitted on all the data (max $950$) gives $1.00$ ---
hiding exactly the fact that the house lies outside the training range.
\end{worked}

\begin{caution}[{{Mistakes that cost marks}}]
\begin{itemize}
  \item \textbf{Sample sd in standardisation.} \texttt{StandardScaler} divides
    by $n$.
  \item \textbf{\texttt{fit\_transform} on the test set.} Fit on training,
    \texttt{transform} the test set.
  \item \textbf{Scaling for a tree} and claiming it helped.
  \item \textbf{Min--max on a column with kept extremes.} Use the robust
    scaler.
\end{itemize}
\end{caution}

\begin{tryit}[{{Exercises 11 --- numerical}}]
\begin{enumerate}[label=\textbf{E11.\arabic*}]
  \item Training mean $80$, sd $20$: standardise $x = 100$. Training min $10$,
    max $30$: min--max scale a held-out $40$ without clipping.
  \item Area ranges from $500$ to $5000$: min--max scale $x = 2000$.
  \item For $\{10, 20, 30, 40, 50\}$ compute $\mu$, $\sigma$ (divisor $n$) and
    the $z$-score of $50$.
  \item Using the data of E10.2 (median $37.5$, $\IQR = 20$), robust-scale
    $x = 60$.
  \item Query (1800 sq ft, 3 bed); A (1850, 1); B (1400, 3). Distances
    unscaled, then with Area $\in [1000, 3000]$ and Bedrooms $\in [1, 5]$
    min--max scaled. Nearest in each case?
  \item Put $v = (-4, 0, 2, 6, 8, 16)$ through all four scalers (hinges for the
    robust scaler).
\end{enumerate}
\end{tryit}

\begin{concept}
\begin{enumerate}[label=\textbf{C11.\arabic*}]
  \item \emph{Normalisation vs standardisation, one point.} Normalisation
    (min--max) maps to a range, $[0,1]$; standardisation maps to mean $0$ and
    sd $1$ and is not bounded.
  \item \emph{Why does a decision tree need no scaling?} A split asks ``is
    $x_j < t$?''. Any increasing rescaling moves $t$ but not which rows fall
    on each side.
  \item \emph{Features range from $0.001$ to $5000$ and you use an RBF SVM.}
    Standardise: the kernel uses squared distance, which the large feature
    would otherwise dominate.
\end{enumerate}
\end{concept}
%======================================================================
\section{Encoding and feature engineering}
\label{sec:encoding}

\method{turning text into numbers, and building better features}
\begin{enumerate}
  \item \textbf{Nominal} (no order: Location) $\to$ one-hot, one $0/1$ column
    per level. With an intercept and interpreted coefficients, drop one level
    (\texttt{drop='first'}) to avoid the dummy-variable trap. Set
    \texttt{handle\_unknown='ignore'} for levels unseen in training.
  \item \textbf{Ordinal} (real order: Condition) $\to$ integer codes \emph{in
    that order}, given explicitly: Poor $0$, Average $1$, Good $2$,
    Excellent $3$.
  \item \textbf{Binary} (Garage) $\to$ one $0/1$ column.
  \item \textbf{Columns after one-hot}: original columns $- 1 + k$ for a
    column with $k$ levels ($-1 + k - 1$ with drop-first).
  \item \textbf{Feature engineering}: create columns from domain knowledge ---
    ratios, sums, differences, date parts --- \emph{only} from information
    available at prediction time. A feature computed from the target is a
    leak.
\end{enumerate}
\end{keyidea}

\begin{worked}[{{Example 12.1 --- encoding the housing table}}]
\textbf{Location} (4 levels) $\to$ \texttt{Loc\_ClementTown},
\texttt{Loc\_PremNagar}, \texttt{Loc\_RajpurRoad}, \texttt{Loc\_Sahastradhara};
H02 (Rajpur Road) becomes $(0, 0, 1, 0)$. With drop-first, 3 columns and
Clement Town is the baseline $(0,0,0)$.

\textbf{Why not codes $0,1,2,3$ for Location?} They impose an order and equal
spacing that do not exist: a linear model must then put Prem Nagar's effect
exactly between Clement Town's and Rajpur Road's, and $k$-NN would call
Clement Town ``closer'' to Prem Nagar than to Sahastradhara. With three
colours coded $0, 1, 2$ and mean targets $20, 50, 10$, the best line through
$(0,20), (1,50), (2,10)$ is $\hat y = 31.67 - 5x$, with fitted values
$31.67, 26.67, 21.67$ and $\mathrm{SSE} = 816.7$ (against $866.7$ for
predicting the overall mean): the codes explain almost nothing. One-hot fits
the three means exactly ($\mathrm{SSE} = 0$). (Trees are fine with integer codes.)

\textbf{Dummy-variable trap.} With an intercept and all $k$ dummies, the
dummies sum to $1$ in every row --- the intercept column --- so $X^\top X$ is
singular and the coefficients are not unique. Drop one level (or use ridge/lasso,
whose penalty makes the solution unique).

\textbf{Unseen level.} A test house in Mussoorie gets $(0,0,0,0)$ under
\texttt{handle\_unknown='ignore'}: it matches no known locality and receives
only the intercept (with drop-first it would be mistaken for the baseline), so
its prediction should be flagged as less trustworthy. Label encoding would have
crashed or silently assigned a meaningless new code.
\end{worked}

\begin{worked}[{{Example 12.2 --- engineering a feature}}]
H09 (1350 sq ft, 3 bed, 1 bath, 46 lakh) and H02 (1500 sq ft, 3 bed,
2 bath, 78 lakh) look nearly identical on Area and Bedrooms, yet H02 costs
$70\%$ more. \texttt{Area\_per\_Bedroom}: $450$ vs $500$;
\texttt{Bath\_per\_Bed}: $0.33$ vs $0.67$ --- the layout information was in
the table, but only in the \emph{relationship} between columns.

For a 2400 sq ft house with 3 bedrooms and 2 bathrooms:
\texttt{Area\_per\_Bedroom} $= 800$, \texttt{Total\_Rooms} $= 5$.
\texttt{Age\_at\_sale} is safe (known when you price the house);
\texttt{Price\_per\_sqft} is not --- it is computed from the price you are
trying to predict.
\end{worked}

\begin{tryit}[{{Exercises 12 --- numerical}}]
\begin{enumerate}[label=\textbf{E12.\arabic*}]
  \item Three known localities, full one-hot: how many indicator columns?
  \item A dataset has 500 rows and 6 columns, one categorical with 5 levels.
    Columns after one-hot encoding it? With drop-first?
  \item Codes blue $0$, green $1$, red $2$; mean targets $30$, $10$, $20$. Fit
    the least-squares line through the three points and give its SSE. What
    would one-hot give?
\end{enumerate}
\end{tryit}

\begin{concept}
\begin{enumerate}[label=\textbf{C12.\arabic*}]
  \item \emph{Blood group (A, B, AB, O) for a linear model? Education
    (School, Bachelor, Master, PhD) for a linear model?} Blood group is
    nominal $\Rightarrow$ one-hot. Education has a real order $\Rightarrow$
    ordinal codes $0$--$3$ in that order.
  \item \emph{When is label encoding nevertheless fine?} For tree-based models,
    which only split on order within one column.
  \item \emph{Encoding vs feature engineering in one line.} Encoding
    translates an existing column into numbers; engineering creates new
    information from existing columns.
  \item \emph{Why is \texttt{OrdinalEncoder}'s default wrong for \{low, medium,
    high\}?} It sorts alphabetically --- high $0$, low $1$, medium $2$ ---
    scrambling the order. Pass the categories explicitly.
\end{enumerate}
\end{concept}

%======================================================================
\section{Data reduction and PCA}
\label{sec:pca}

\method{reducing columns, and PCA by hand}
\begin{enumerate}
  \item \textbf{Three kinds of reduction}: \emph{feature selection} keeps a
    subset of the original columns (filter, wrapper, embedded);
    \emph{dimensionality reduction} (feature extraction, e.g.\ PCA) builds new
    columns from all of them; \emph{sampling} (numerosity reduction) keeps
    fewer rows. Selection cuts columns; sampling cuts rows.
  \item \textbf{PCA}: standardise if the units differ; centre; covariance
    matrix $S$; solve $\det(S - \lambda I) = 0$ (for $2\times2$,
    $\lambda^2 - \operatorname{tr}(S)\lambda + \det S = 0$); unit eigenvectors
    from $(S - \lambda I)u = 0$.
  \item \textbf{Shortcut} for two standardised columns with correlation $r$:
    $S = \begin{pmatrix}1 & r\\ r & 1\end{pmatrix}$, $\lambda = 1 \pm r$,
    $u = \frac{1}{\sqrt2}(1, \pm1)$.
  \item Explained variance $\lambda_k/\sum\lambda$; scores $t = x_c \cdot u$.
    Check $\sum\lambda = \operatorname{tr}S$, $\prod\lambda = \det S$.
\end{enumerate}
\end{keyidea}

\begin{worked}[{{Example 13.1 --- five houses, Area and Bedrooms}}]
Area $(1000, 1500, 2000, 2500, 3000)$, Bedrooms $(2, 3, 4, 4, 5)$.
\textbf{Standardise} (divisor $n$): Area $\mu = 2000$, $\sigma = 707.11$,
$z_A = (-1.4142, -0.7071, 0, 0.7071, 1.4142)$; Bedrooms $\mu = 3.6$,
$\sigma = \sqrt{5.2/5} = 1.0198$, $z_B = (-1.5689, -0.5883, 0.3922, 0.3922,
1.3728)$.

\textbf{Correlation} $r = \frac15\sum z_Az_B = 4.8536/5 = 0.9707$, so
$S = \begin{pmatrix}1 & 0.9707\\0.9707 & 1\end{pmatrix}$.
\textbf{Eigenvalues} $\lambda_1 = 1.9707$, $\lambda_2 = 0.0293$ (sum $2$);
PC1 explains $1.9707/2 = 98.5\%$. \textbf{Directions}
$u_1 = \frac1{\sqrt2}(1,1)$ (``overall size''), $u_2 = \frac1{\sqrt2}(1,-1)$.
\textbf{Scores} PC1 $= (z_A + z_B)/\sqrt2$: H01 $(-1.4142 - 1.5689)/1.4142 =
-2.1094$; all five: $-2.1094, -0.9160, 0.2774, 0.7774, 1.9707$. Two columns
replaced by one, $1.5\%$ of the variance lost --- because the two columns
said nearly the same thing.
\end{worked}

\begin{worked}[{{Example 13.2 --- a covariance matrix that needs the full method}}]
Points $(1,1), (3,3), (4,3), (5,5), (7,3)$; mean $(4,3)$; centred
$(-3,-2), (-1,0), (0,0), (1,2), (3,0)$. Sums $\sum x_1^2 = 20$,
$\sum x_2^2 = 8$, $\sum x_1x_2 = 8$; divide by $n - 1 = 4$:
$S = \begin{pmatrix}5 & 2\\2 & 2\end{pmatrix}$.
$\lambda^2 - 7\lambda + 6 = 0 \Rightarrow \lambda = 6, 1$ (check: $6+1 = 7$,
$6 \times 1 = 6$). $\lambda = 6$: $-u_1 + 2u_2 = 0 \Rightarrow u =
(2,1)/\sqrt5$; $\lambda = 1$: $u = (-1,2)/\sqrt5$. PC1 explains $6/7 =
85.7\%$. Scores $(2x_{c1} + x_{c2})/\sqrt5 = (-8,-2,0,4,6)/\sqrt5$, sample
variance $120/20 = 6 = \lambda_1$.
\end{worked}

\begin{caution}[{{Mistakes that cost marks}}]
\begin{itemize}
  \item \textbf{Not scaling first} --- PCA then finds the column with the
    largest numbers.
  \item \textbf{Calling PCA feature selection.} Every component mixes every
    original feature; ``Age was not useful'' cannot be concluded from PCA.
  \item \textbf{Forgetting PCA is unsupervised}: it never looks at $y$, and
    the discarded low-variance direction may be the predictive one.
  \item \textbf{Unnormalised eigenvectors} in the scores.
\end{itemize}
\end{caution}

\begin{tryit}[{{Exercises 13 --- numerical}}]
\begin{enumerate}[label=\textbf{E13.\arabic*}]
  \item Unit principal direction $u = (1/\sqrt2, 1/\sqrt2)$; project the centred
    point $(2, 2)$ onto it.
  \item Two standardised features have $r = 0.8$. Write the covariance matrix,
    both eigenvalues, and the variance explained by PC1.
  \item Points $(1,3), (5,3), (5,4), (7,3), (7,7)$: $S$ (divisor $n-1$),
    eigenvalues, unit eigenvectors, PC1's share, and the PC1 scores.
  \item Eigenvalues $5.2, 2.1, 1.1, 0.9, 0.4, 0.3$. How many components keep
    $90\%$ of the variance? $95\%$?
\end{enumerate}
\end{tryit}

\begin{concept}
\begin{enumerate}[label=\textbf{C13.\arabic*}]
  \item \emph{Does PCA keep a subset of the named columns?} No --- it builds
    new columns, each a weighted mix of all of them.
  \item \emph{50 features: PCA keeps 5 components (92\% variance); a colleague
    selects 5 original features. Which can you explain to a manager? Which
    helps $k$-NN more?} Selection --- the features keep their meaning and
    units. PCA usually helps $k$-NN more: its components are uncorrelated and
    carry information from all 50 columns.
  \item \emph{Why might keeping 95\% of the variance hurt prediction? How
    should $k$ be chosen?} Variance in $X$ is not relevance to $y$. Choose the
    number of components by cross-validated prediction error, inside the
    pipeline.
  \item \emph{Curse of dimensionality: what and whom?} As dimensions grow the
    space's volume grows exponentially, data become sparse and all points
    nearly equidistant. $k$-NN, $k$-means and RBF SVMs suffer directly.
  \item \emph{3000 features, 200 rows?} Reduce (PCA or lasso) or regularise
    strongly, selecting inside cross-validation; a linear model could fit the
    training data perfectly.
\end{enumerate}
\end{concept}

%======================================================================
\section{Splitting and cross-validation}
\label{sec:splits}

\method{honest evaluation}
\begin{enumerate}
  \item \textbf{Train / validation / test}: fit on training; choose
    hyperparameters on validation; report \emph{once} on test.
  \item \textbf{Split counts}: $n$ rows at ratio $a{:}b$ $\to$ $na/(a+b)$ and
    $nb/(a+b)$.
  \item \textbf{$k$-fold CV}: $k$ folds of $n/k$ rows; each is the
    validation fold once; $k$ fits, each on $n(1 - 1/k)$ rows. Report the
    mean \emph{and} the spread of the $k$ scores. LOOCV is $k = n$.
  \item \textbf{Match the split to the data}: \emph{stratified} (keep class
    proportions in every fold), \emph{grouped} (all rows of one patient or
    customer on one side), \emph{time-ordered} (train on the past, test on the
    future).
  \item A random split of $n$ rows with $n_+$ positives puts $k$ positives in
    a test set of $t$ rows with probability
    $\binom{n_+}{k}\binom{n - n_+}{t - k}/\binom nt$.
\end{enumerate}
\end{keyidea}

\begin{worked}[{{Example 14.1 --- counts, and the mean and spread of CV}}]
$400$ houses split $75{:}25$: $300$ training, $100$ test. 5-fold CV on $100$
rows: $20$ validation rows per fold, $5$ fits of $80$ rows each.

Fold scores $0.78, 0.82, 0.75, 0.80, 0.85$: mean $4.00/5 = 0.80$; deviations
$-0.02, 0.02, -0.05, 0, 0.05$; squares sum to $0.0058$; sd $=
\sqrt{0.0058/5} = 0.0341$ (divisor $k$; with $k - 1$, $0.0381$ --- say which).
Two models whose CV means differ by less than about this spread have not been
shown to differ.
\end{worked}

\begin{worked}[{{Example 14.2 --- why stratify, group and respect time}}]
\textbf{Twenty patients, five ill, four held out.} At random, the chance of
\emph{no} ill patient in the test set is
$\binom{15}{4}/\binom{20}{4} = 1365/4845 = 0.2817$ --- more than one split in
four leaves recall uncomputable. Stratified: $5 \times \tfrac{4}{20} = 1$ ill
and $3$ healthy, every time.

\textbf{Grouped.} Repeated measurements of the same patient (or many
transactions per customer) in both train and test let the model recognise
individuals: the score is inflated. Split by patient (\texttt{GroupKFold}).
\textbf{Time.} Stock prices split at random let the model train on the future
and interpolate; use a temporal split.
\end{worked}

\begin{tryit}[{{Exercises 14 --- numerical}}]
\begin{enumerate}[label=\textbf{E14.\arabic*}]
  \item $400$ rows split $75{:}25$: rows in the training portion?
  \item 5-fold CV on $100$ rows: validation rows per fold, and model fits?
  \item Fold scores $0.81, 0.79, 0.86, 0.74, 0.80$: mean and sd (divisor $k$).
  \item $30$ rows, $6$ positive, $5$ held out at random: probability of no
    positive in the test set? What does a stratified split give?
\end{enumerate}
\end{tryit}

\begin{concept}
\begin{enumerate}[label=\textbf{C14.\arabic*}]
  \item \emph{Use of training, validation and test sets?} Training: learn the
    parameters. Validation: compare models and tune hyperparameters. Test: a
    single, final, unbiased estimate of performance on new data.
  \item \emph{``My model is 96\% accurate'' --- measured on training data.
    What is wrong, and what instead?} A training score measures fit, not
    generalisation, and rewards memorisation. Evaluate on a held-out test set
    (or by cross-validation).
  \item \emph{Why a validation set when a test set exists?} Choosing a
    hyperparameter by the test score is fitting to the test set; after many
    comparisons the winner partly reflects that set's noise. The validation
    set absorbs this so the test set stays clean.
  \item \emph{Only 60 observations --- how to estimate performance?} $k$-fold
    CV (or LOOCV); a single test set of a few rows is too unstable.
  \item \emph{Tuning 50 models on the same CV scores and reporting the best?}
    Selection optimism: the best of 50 noisy scores is biased upward. Use
    nested CV, or confirm on an untouched test set.
\end{enumerate}
\end{concept}

%======================================================================
\section{Imbalanced data}
\label{sec:imbalance}

\method{metrics and remedies for a rare class}
\begin{enumerate}
  \item Build the confusion matrix for the rare (positive) class: $TP$,
    $FP$, $FN$, $TN$.
  \item accuracy $= \frac{TP + TN}{\text{all}}$; precision $= \frac{TP}{TP +
    FP}$ (of those flagged, how many were real); recall $= \frac{TP}{TP + FN}$
    (of the real ones, how many were caught); $F_1 = \frac{2PR}{P + R}$.
  \item Compare with the \textbf{always-negative} baseline: its accuracy
    equals the majority share, its recall is $0$.
  \item \textbf{Remedies}, applied to the training part only: random
    oversampling (duplicate minority rows), random undersampling (discard
    majority rows), SMOTE ($x_{\text{new}} = x_i + \lambda(x_{\text{nb}} -
    x_i)$, $\lambda \in [0,1]$), class weights, a moved threshold.
\end{enumerate}
\end{keyidea}

\begin{worked}[{{Example 15.1 --- one thousand patients, fifty ill}}]
\textbf{Always ``no disease''}: $TN = 950$, $FN = 50$, $TP = FP = 0$. Accuracy
$950/1000 = 95\%$; recall $0/50 = 0$; precision $0/0$, undefined.

\textbf{A real model} flags 100 patients, 40 truly ill: $TP = 40$, $FP = 60$,
$FN = 10$, $TN = 890$.
\[
  \text{precision} = \tfrac{40}{100} = 0.40, \quad
  \text{recall} = \tfrac{40}{50} = 0.80, \quad
  F_1 = \tfrac{2(0.40)(0.80)}{1.20} = 0.533,
\]
\[
  \text{accuracy} = \tfrac{930}{1000} = 93\% .
\]
Lower accuracy than the useless model, yet it catches $80\%$ of the sick. In
screening that trade is worth making; whether to favour precision or recall
is a decision about consequences, not statistics.
\end{worked}

\begin{worked}[{{Example 15.2 --- resampling arithmetic, and SMOTE}}]
From $950 + 50$: random oversampling $\to 950 + 950 = 1900$ rows (the 50 ill
patients repeated); random undersampling $\to 50 + 50 = 100$ rows ($900$
patients, $90\%$ of the data, thrown away). Balanced class weights
$n/(K n_c)$: $1000/1900 = 0.5263$ and $1000/100 = 10$.

\textbf{SMOTE} between $x_i = (5, 120)$ and its neighbour $(8, 150)$ with
$\lambda = 0.25$: $x_i + 0.25(3, 30) = (5.75, 127.5)$. Every synthetic point
lies on the segment between two real minority points --- unsafe when the
classes interleave, with categorical features, or in high dimensions.
\end{worked}

\begin{caution}[{{Mistakes that cost marks}}]
\begin{itemize}
  \item \textbf{Reporting accuracy} on an imbalanced problem.
  \item \textbf{Resampling before the split}, or resampling the test set:
    copies or interpolations of test rows leak into training, and a resampled
    test set describes a world with the wrong prevalence.
  \item \textbf{Swapping precision and recall.}
\end{itemize}
\end{caution}

\begin{tryit}[{{Exercises 15 --- numerical}}]
\begin{enumerate}[label=\textbf{E15.\arabic*}]
  \item Among $400$ rows, $40$ are positive: the positive percentage? The
    accuracy of an always-negative classifier?
  \item Of $2000$ transactions, $40$ are fraudulent. A model flags $90$ and
    catches $30$ of the frauds. Confusion matrix, accuracy, precision, recall,
    $F_1$.
  \item $10{,}000$ transactions, $100$ fraudulent; a detector flags $150$ and
    catches $60$. Confusion matrix and all four metrics. A rival flags
    nothing: its accuracy? Which would you deploy?
  \item $1200$ majority and $80$ minority rows: rows after random oversampling
    to balance? After undersampling? How many discarded?
  \item $10{,}000$ patients, $200$ with a condition: accuracy and recall of
    ``always healthy''?
\end{enumerate}
\end{tryit}

\begin{concept}
\begin{enumerate}[label=\textbf{C15.\arabic*}]
  \item \emph{400 houses, only 40 expensive: one sampling method; does it add
    or remove rows?} Random oversampling (or SMOTE) adds minority rows; random
    undersampling removes majority rows.
  \item \emph{One application where recall matters more, one for precision.}
    Recall: disease screening (a miss is costly). Precision: automatically
    suspending accounts (a false alarm is costly).
  \item \emph{Oversampling vs SMOTE.} Oversampling duplicates existing rows
    (the model can memorise them); SMOTE interpolates new points (it asserts
    that the segment between two minority points is minority territory).
\end{enumerate}
\end{concept}

%======================================================================
\section{Data leakage and the order of the pipeline}
\label{sec:leakage}

\method{finding a leak}
\begin{enumerate}
  \item \textbf{Target leakage}: a feature that would not exist at the moment
    the prediction is made, or is computed from the target. Test: \emph{at
    prediction time, would I have this value?}
  \item \textbf{Preprocessing leakage}: any statistic (mean, median, min/max,
    quartiles, selected features, PCA axes, resampling) learned from rows
    that include the test set.
  \item \textbf{The safe order}: understand the columns $\to$ drop leaks and
    identifiers $\to$ correct known errors $\to$ \textbf{split} $\to$ impute,
    transform, encode, scale, select, resample --- all \emph{fitted on
    training data} inside a \texttt{Pipeline} $\to$ model $\to$ evaluate once
    on test.
  \item A pipeline prevents preprocessing leakage; only domain judgement
    prevents target leakage.
\end{enumerate}
\end{keyidea}

\begin{worked}[{{Example 16.1 --- find every leak}}]
\begin{lstlisting}
df["price_per_sqft"] = df["Price_lakh"] / df["Area_sqft"]
df = df.fillna(df.mean(numeric_only=True))
X = StandardScaler().fit_transform(df.drop(columns=["Price_lakh"]))
X = SelectKBest(k=5).fit_transform(X, df["Price_lakh"])
\end{lstlisting}
\begin{enumerate}[label=(\roman*)]
  \item \texttt{price\_per\_sqft} is computed from the target: \textbf{target
    leakage} --- the model is handed the answer.
  \item \texttt{fillna(df.mean())} on all rows: \textbf{preprocessing
    leakage} (and the mean is a poor fill for skewed columns).
  \item \texttt{StandardScaler().fit\_transform} on all rows:
    \textbf{preprocessing leakage}.
  \item \texttt{SelectKBest} with the labels of all rows:
    \textbf{preprocessing leakage} through feature selection --- and it
    survives later cross-validation, because every fold inherits the same
    contaminated columns.
\end{enumerate}
Also in the housing table: \textbf{Buyer\_Income} is known only after the
sale --- a target leak; \textbf{ID} is an identifier --- drop it (data
reduction: feature selection).
\end{worked}

\begin{worked}[{{Example 16.2 --- the pipeline, in order}}]
\small
\begin{enumerate}
  \item Drop \texttt{ID} and \texttt{Income}; correct H12's Age $150 \to 15$.
  \item Split into training and test sets (stratify or group if needed).
  \item On the training set only: median-impute Area and Age, group-wise
    mode for Garage; log-transform Price (or keep H11 and use a robust
    method); one-hot Location, ordinal Condition, binary Garage; engineer
    \texttt{Area\_per\_Bedroom}; scale if the model needs it.
  \item Wrap step 3 and the model in a \texttt{Pipeline}, tune by
    cross-validation on the training set, then evaluate \emph{once} on the
    test set.
\end{enumerate}
\end{worked}

\begin{concept}
\begin{enumerate}[label=\textbf{C16.\arabic*}]
  \item \emph{\texttt{scaler.fit\_transform(X\_test)} appears in a pipeline.}
    It learns new parameters from the test data, so train and test are scaled
    differently and the test set influences the transformation. Use
    \texttt{fit\_transform} on training, \texttt{transform} on test.
  \item \emph{A readmission model uses \texttt{discharge\_destination}.} Target
    leakage for any model meant to predict at admission: the value is recorded
    at discharge and reflects the outcome being predicted.
  \item \emph{SMOTE on the whole dataset, then split, then 99\% accuracy.}
    Synthetic points derived from rows that later land in the test set sit in
    training --- the model has seen smeared copies of the answers.
  \item \emph{The 50 features most correlated with the target are selected,
    then the data are split and cross-validated; scores collapse in
    production.} Selection leakage: the test rows helped choose the features,
    and CV cannot catch it because selection happened before the loop. Put
    selection inside the pipeline.
  \item \emph{A house-price model uses the final registered sale tax.} Target
    leakage --- the tax is computed from the sale price. Remove it and
    evaluate on data that mimics prediction time.
  \item \emph{Loan default (3\% default): which column goes first ---
    \texttt{days\_until\_first\_missed\_payment} or
    \texttt{loan\_officer\_id}?} The first: it exists only after a default
    has begun --- a textbook target leak. The officer ID is an identifier and
    should also go.
  \item \emph{Why is leakage more dangerous than a crash?} It makes results
    look \emph{better}, so every incentive is to accept it.
  \item \emph{Why must SMOTE, scaling and selection sit inside the CV loop?}
    So that every statistic is refitted on each fold's training part and the
    held-out fold stays unseen.
\end{enumerate}
\end{concept}
\unit{Unit V --- Regression (L8--L9)}
%======================================================================
\section{Regression foundations}
\label{sec:regfound}

\method{setting up a regression problem}
\begin{enumerate}
  \item \textbf{Regression} predicts a continuous numerical target from one
    or more features, by estimating $f$ with $y \approx f(x)$ --- for the data
    given and for new data.
  \item \textbf{Two jobs}: \emph{estimation} (how much does a square foot add?)
    and \emph{prediction} (what will this 1750 sq ft house fetch?).
  \item \textbf{``Linear'' means linear in the parameters}, not in $x$:
    $b_0 + b_1x + b_2x^2$ and $b_0 + b_1\ln x$ are linear models;
    $b_0e^{b_1x}$ is not.
  \item \textbf{Eight steps}: state the question $\to$ collect and plot $\to$
    choose the form $\to$ split $\to$ estimate (least squares) $\to$ check
    adequacy (residuals) $\to$ validate $\to$ interpret and report. If the
    residuals say the form is wrong, go back to step 3.
  \item \textbf{Simple linear regression}: $y_i = \beta_0 + \beta_1x_i +
    \varepsilon_i$; fitted $\hat y_i = b_0 + b_1x_i$. Greek $=$ unknown
    population values; Latin $=$ estimates from the sample. $\varepsilon_i$
    (from the true line) is unobservable; the residual $e_i = y_i - \hat y_i$
    is computed.
  \item \textbf{Assumptions}, most to least damaging when violated:
    linearity, independence, homoscedasticity (constant $\operatorname{Var}
    \varepsilon$), normality.
\end{enumerate}
\end{keyidea}

\begin{worked}[{{Example 17.1 --- regression or not, linear or not}}]
\small
\begin{tabular}{@{}p{5.6cm}p{8.3cm}@{}}
\toprule
\textbf{Model / task} & \textbf{Verdict}\\ \midrule
Predict tomorrow's closing price, yield per hectare, length of hospital stay
& regression: the target is continuous\\
Predict whether a house is ``expensive'' & classification: the target is a
category\\
$\hat y = b_0 + b_1x + b_2x^2$ & linear regression: make a column $x^2$ and it is
a two-feature linear model\\
$\hat y = b_0 + b_1\ln x$ & linear regression in the feature $\ln x$\\
$\hat y = b_0e^{b_1x}$ & \emph{not} linear: $b_1$ sits inside the exponential
(taking $\ln y$ turns it into a linear model of $\ln y$)\\
\bottomrule
\end{tabular}

\normalsize
\textbf{$\beta$ vs $b$, $\varepsilon$ vs $e$.} A different sample of 200 houses
gives different $b$'s and the same $\beta$'s. The error term carries
everything not in the model --- condition of the plaster, the seller's
urgency --- and the model asks only that it be unsystematic.
\end{worked}

\begin{concept}
\begin{enumerate}[label=\textbf{C17.\arabic*}]
  \item \emph{State the simple linear regression model and name every
    symbol.} $y_i = \beta_0 + \beta_1x_i + \varepsilon_i$: $y$ target, $x$
    feature, $\beta_0$ intercept, $\beta_1$ slope, $\varepsilon_i$ unobserved
    error, $i = 1,\dots,n$.
  \item \emph{Why squared error to derive the estimates? Two reasons.} It is
    differentiable and convex, giving a unique closed-form solution; and it
    punishes large misses more than small ones, which is usually wanted (and,
    by the Gauss--Markov theorem, gives the minimum-variance linear unbiased
    estimator). The price: one outlier can drag the line.
  \item \emph{Name the four assumptions and rank them.} Linearity,
    independence, homoscedasticity, normality --- in that order of damage.
    Heteroscedasticity leaves the coefficients unbiased but the intervals
    wrong; normality matters only for small-sample inference.
  \item \emph{Estimation vs prediction?} Estimation asks for a trustworthy
    coefficient; prediction asks for a good $\hat y$ on new data. A model can
    do one well and the other badly.
\end{enumerate}
\end{concept}

%======================================================================
\section{Least squares by hand}
\label{sec:ls}

\method{fitting $\hat y = b_0 + b_1x$}
\begin{enumerate}
  \item Minimise $S(b_0, b_1) = \sum_i (y_i - b_0 - b_1x_i)^2$.
  \item \textbf{Normal equations} (set both partial derivatives to zero):
    $nb_0 + b_1\sum x_i = \sum y_i$ and $b_0\sum x_i + b_1\sum x_i^2 = \sum
    x_iy_i$.
  \item Solution: $b_1 = S_{xy}/S_{xx}$, $b_0 = \bar y - b_1\bar x$, with
    $S_{xx} = \sum(x_i - \bar x)^2 = \sum x_i^2 - n\bar x^2$ and
    $S_{xy} = \sum(x_i - \bar x)(y_i - \bar y) = \sum x_iy_i - n\bar x\bar y$.
  \item \textbf{Table}: $x - \bar x$, $y - \bar y$, $(x - \bar x)^2$,
    $(x - \bar x)(y - \bar y)$, then $\hat y$, $e$, $e^2$.
  \item \textbf{Checks} (they are the two normal equations): $\sum e_i = 0$,
    $\sum x_ie_i = 0$; the line passes through $(\bar x, \bar y)$;
    $\mathrm{SST} = \mathrm{SSR} + \mathrm{SSE}$ with $\mathrm{SSR} =
    b_1^2S_{xx}$.
  \item \textbf{Matrix form}: $b = (X^\top X)^{-1}X^\top y$ with a column of
    ones; $X^\top X = \begin{pmatrix} n & \sum x\\ \sum x & \sum
    x^2\end{pmatrix}$, $X^\top y = \begin{pmatrix}\sum y\\ \sum
    xy\end{pmatrix}$.
\end{enumerate}
\end{keyidea}

\begin{worked}[{{Example 18.1 --- derive $b_0$ and $b_1$ (Class Test Q4)}}]
\textbf{What we minimise:} $S(b_0, b_1) = \sum_{i=1}^{n}(y_i - b_0 -
b_1x_i)^2$.

\textbf{Differentiate.}
\[
  \frac{\partial S}{\partial b_0} = -2\sum(y_i - b_0 - b_1x_i) = 0,
  \qquad
  \frac{\partial S}{\partial b_1} = -2\sum x_i(y_i - b_0 - b_1x_i) = 0 .
\]
\textbf{Normal equations.} $nb_0 + b_1\sum x_i = \sum y_i$ \ (1), and
$b_0\sum x_i + b_1\sum x_i^2 = \sum x_iy_i$ \ (2).

\textbf{Solve.} Divide (1) by $n$: $b_0 + b_1\bar x = \bar y$, so $b_0 =
\bar y - b_1\bar x$. Substitute into (2) and use $\sum x_i = n\bar x$:
\[
  b_1\Bigl(\sum x_i^2 - n\bar x^2\Bigr) = \sum x_iy_i - n\bar x\bar y
  \quad\Longrightarrow\quad
  b_1 = \frac{\sum x_iy_i - n\bar x\bar y}{\sum x_i^2 - n\bar x^2}
      = \frac{S_{xy}}{S_{xx}} .
\]
Defined only when $S_{xx} > 0$ (the $x$ values are not all equal). The
second derivatives ($2n > 0$ and a positive determinant $4nS_{xx}$) confirm a
minimum.

\textbf{The line passes through $(\bar x, \bar y)$}: at $x = \bar x$,
$\hat y = b_0 + b_1\bar x = (\bar y - b_1\bar x) + b_1\bar x = \bar y$.
\end{worked}

\begin{worked}[{{Example 18.2 --- six houses by hand}}]
\begin{center}
\small
\begin{tabular}{@{}lrrrrrr|r@{}}
\toprule
house & A & B & C & D & E & F & sum\\ \midrule
Area $x$ & 600 & 800 & 1000 & 1200 & 1400 & 1600 & 6600\\
Price $y$ & 40 & 61 & 74 & 81 & 87 & 107 & 450\\ \midrule
$x - \bar x$ & $-500$ & $-300$ & $-100$ & $100$ & $300$ & $500$ & $0$\\
$y - \bar y$ & $-35$ & $-14$ & $-1$ & $6$ & $12$ & $32$ & $0$\\
$(x - \bar x)^2$ & 250000 & 90000 & 10000 & 10000 & 90000 & 250000 & \textbf{700000}\\
$(x - \bar x)(y - \bar y)$ & 17500 & 4200 & 100 & 600 & 3600 & 16000 & \textbf{42000}\\ \midrule
$\hat y = 9 + 0.06x$ & 45 & 57 & 69 & 81 & 93 & 105 & 450\\
$e = y - \hat y$ & $-5$ & $+4$ & $+5$ & $0$ & $-6$ & $+2$ & $0$\\
$e^2$ & 25 & 16 & 25 & 0 & 36 & 4 & \textbf{106}\\
\bottomrule
\end{tabular}
\end{center}
$\bar x = 1100$, $\bar y = 75$; $b_1 = 42000/700000 = 0.06$ lakh per sq ft;
$b_0 = 75 - 0.06 \times 1100 = 9$:
\[
  \boxed{\hat y = 9 + 0.06\,x}
\]
\textbf{Checks.} $\sum e = 0$; $\sum xe = 0$; at $x = 1100$, $\hat y = 75$.
$\mathrm{SSE} = 106$, $\mathrm{MSE} = 106/6 = 17.67$. $\mathrm{SST} = S_{yy} =
2626$, $\mathrm{SSR} = 2626 - 106 = 2520 = 0.06^2 \times 700000$ \checkmark.
No other straight line has a smaller SSE: $S(b_0, b_1)$ is a convex bowl with
one bottom at $(9, 0.06)$ --- which is why calculus, not gradient descent, is
enough here. (Gradient descent on the same SSE reaches the same answer, more slowly --- and much faster when the features are scaled.)

\textbf{Interpretation.} Each extra square foot is \emph{associated with}
$0.06$ lakh (Rs.~6{,}000) more;
$b_0 = 9$ is the value at $x = 0$, far outside the data and not meaningful on
its own.
\end{worked}

\begin{worked}[{{Example 18.3 --- the matrix route}}]
$x = (0, 1, 2, 3)$, $y = (1, 3, 4, 6)$:
$X^\top X = \begin{pmatrix}4 & 6\\ 6 & 14\end{pmatrix}$,
$X^\top y = \begin{pmatrix}14\\ 29\end{pmatrix}$, $\det = 56 - 36 = 20$,
\[
  b = \frac1{20}\begin{pmatrix}14 & -6\\ -6 & 4\end{pmatrix}
      \begin{pmatrix}14\\29\end{pmatrix}
    = \frac1{20}\begin{pmatrix}196 - 174\\ -84 + 116\end{pmatrix}
    = \begin{pmatrix}1.1\\ 1.6\end{pmatrix}.
\]
Cross-check: $\bar x = 1.5$, $\bar y = 3.5$, $S_{xx} = 5$, $S_{xy} = 8$,
$b_1 = 1.6$, $b_0 = 1.1$ \checkmark.

\textbf{How precise is $b_1$?} $\operatorname{Var}(b_1) = \sigma^2/S_{xx}$:
spreading the $x$ values out (larger $S_{xx}$) pins the slope down more
precisely, at no extra cost in $n$.
\end{worked}

\begin{caution}[{{Mistakes that cost marks}}]
\begin{itemize}
  \item \textbf{$b_1 = S_{xy}/S_{yy}$} --- that is the slope of $x$ on $y$.
  \item \textbf{Rounding $b_1$ before computing $b_0$.}
  \item \textbf{Residuals that do not sum to zero} mean an arithmetic error,
    every time (with an intercept).
  \item \textbf{Extrapolating} far outside the range of $x$, or reading $b_1$
    causally.
\end{itemize}
\end{caution}

\begin{tryit}[{{Exercises 18 --- numerical}}]
\begin{enumerate}[label=\textbf{E18.\arabic*}]
  \item $y = 12$, $\hat y = 10$: the squared residual. For six houses with
    $\mathrm{SSE} = 106$: the MSE.
  \item Least-squares slope and intercept for $(1,2), (2,3), (3,5)$.
  \item Use $\hat y = 9 + 0.06x$ at $x = 600$ with observed $y = 40$: the signed
    residual.
  \item $x = (2, 4, 6, 8, 10)$, $y = (5, 11, 14, 20, 25)$: $\bar x$, $\bar y$,
    $S_{xx}$, $S_{xy}$, $b_1$, $b_0$; predict at $x = 7$; verify the residuals
    sum to zero and compute SSE.
  \item From summary statistics: $n = 10$, $\sum x = 50$, $\sum y = 120$,
    $\sum x^2 = 310$, $\sum xy = 660$. Find $b_1$, $b_0$ and $\hat y(7)$.
  \item Fit $\hat y = wx$ (no intercept) to $x = (1, 2, 3)$, $y = (4, 5, 7)$
    with $w = \sum xy/\sum x^2$. Do the residuals sum to zero? Which property
    survives, and why?
\end{enumerate}
\end{tryit}

\begin{concept}
\begin{enumerate}[label=\textbf{C18.\arabic*}]
  \item \emph{Write the two normal equations and say what each forces on the
    residuals.} $\sum(y_i - b_0 - b_1x_i) = 0$ forces $\sum e_i = 0$;
    $\sum x_i(y_i - b_0 - b_1x_i) = 0$ forces $\sum x_ie_i = 0$ (residuals
    uncorrelated with $x$).
  \item \emph{Why does changing area from sq ft to sq m change $b_1$ but not
    the correlation?} $b_1 = r\,s_y/s_x$: $r$ is unit-free, the slope carries
    the units (a factor of $10.76$ here).
  \item \emph{Model A: train RMSE 2, validation RMSE 9. Model B: 4 and 5.
    Which?} B: A overfits --- its training error flatters it.
\end{enumerate}
\end{concept}
%======================================================================
\section{Direct regression and maximum likelihood}
\label{sec:mle}

\method{choosing ``closest'', and the likelihood route to least squares}
\begin{enumerate}
  \item \textbf{Three defensible lines}, all through $(\bar x, \bar y)$:
    \emph{direct} (vertical distances, $y$ on $x$) slope $S_{xy}/S_{xx}$;
    \emph{inverse} ($x$ on $y$, drawn on $y$--$x$ axes) slope $S_{yy}/S_{xy}$;
    \emph{orthogonal} (perpendicular distances) slope $=$ direction of the
    leading eigenvector of the covariance matrix. Their ratio:
    $\text{direct}/\text{inverse} = r^2$; for $r > 0$ the orthogonal slope lies
    between the other two.
  \item \textbf{Likelihood.} Assume $\varepsilon_i \sim N(0, \sigma^2)$
    independent, so $y_i \mid x_i \sim N(\beta_0 + \beta_1x_i, \sigma^2)$ with
    density $f(y_i) = \frac{1}{\sqrt{2\pi\sigma^2}}\exp\!\bigl(-\frac{(y_i -
    \beta_0 - \beta_1x_i)^2}{2\sigma^2}\bigr)$.
  \item $L = \prod_i f(y_i)$; take logs:
    $\ln L = -\frac n2\ln(2\pi) - \frac n2\ln\sigma^2 - \frac{1}{2\sigma^2}
    \sum(y_i - \beta_0 - \beta_1x_i)^2$.
  \item \textbf{Punchline}: for any fixed $\sigma^2$, maximising $\ln L$ over
    the coefficients $=$ minimising SSE. Under normal errors, least squares
    \emph{is} maximum likelihood.
  \item $\partial\ln L/\partial\sigma^2 = 0$ gives $\hat\sigma^2_{\text{ML}} =
    \mathrm{SSE}/n$ --- biased low. The unbiased estimate is $s^2 =
    \mathrm{SSE}/(n - p - 1)$ ($= \mathrm{SSE}/(n-2)$ for one feature); $s$ is
    the typical residual size, in the units of $y$.
\end{enumerate}
\end{keyidea}

\begin{worked}[{{Example 19.1 --- three lines through the same five points}}]
$x = (1,2,3,4,5)$, $y = (2,4,5,4,5)$: $\bar x = 3$, $\bar y = 4$,
$S_{xx} = 10$, $S_{yy} = 6$, $S_{xy} = 6$.

\textbf{Direct}: $b_1 = 6/10 = 0.6$, $b_0 = 4 - 1.8 = 2.2$.
\textbf{Inverse}: slope $S_{yy}/S_{xy} = 1.0$, intercept $4 - 3 = 1.0$.
\textbf{Orthogonal}: covariance matrix ($n-1 = 4$)
$\begin{pmatrix}2.5 & 1.5\\1.5 & 1.5\end{pmatrix}$, $\lambda =
(4 \pm \sqrt{10})/2 = 3.5811, 0.4189$; leading eigenvector
$(2.5 - 3.5811)v_1 + 1.5v_2 = 0 \Rightarrow$ slope $v_2/v_1 = 0.7208$.

Ratio $0.6/1.0 = 0.6 = r^2 = 36/60$ \checkmark, and $0.6 < 0.7208 < 1.0$. Each
line minimises its own criterion and no other. ``Regression'' means the direct
line; the next example shows why that criterion in particular.
\textbf{Do not use the direct fit backwards}: rearranging $\hat y = 2.2 + 0.6x$
to predict $x$ gives slope $1/0.6$, not the inverse regression.
\end{worked}

\begin{worked}[{{Example 19.2 --- from the likelihood to least squares, on six houses}}]
In $\ln L$ the first two terms do not contain $\beta_0, \beta_1$, and the third
is $-\mathrm{SSE}$ times the positive constant $1/(2\sigma^2)$. So the
$(\beta_0, \beta_1)$ that maximise $\ln L$ are exactly the least-squares
$(b_0, b_1)$: for the six houses, $(9, 0.06)$.

Then
$\frac{\partial \ln L}{\partial\sigma^2} = -\frac{n}{2\sigma^2} +
\frac{\mathrm{SSE}}{2\sigma^4} = 0 \Rightarrow \hat\sigma^2_{\text{ML}} =
\frac{\mathrm{SSE}}{n} = \frac{106}{6} = 17.67$, while
$s^2 = \frac{106}{6 - 2} = 26.5$, $s = 5.15$ lakh. The ratio $(n-2)/n = 4/6$:
at $n = 6$ the ML estimate is a third too small in expectation; at $n = 200$
the gap nearly vanishes.

\textbf{Likelihood is a product of densities.} Two independent observations
with densities $0.2$ and $0.3$ under a candidate model have likelihood
$0.2 \times 0.3 = 0.06$ ($\ln L = -2.8134$). A density height is not the
probability of that exact value --- for a continuous outcome that probability
is $0$.
\end{worked}

\begin{caution}[{{Mistakes that cost marks}}]
\begin{itemize}
  \item \textbf{Dividing SSE by $n$ and calling it unbiased.} That is the ML
    estimate; the unbiased one divides by $n - 2$ (one feature).
  \item \textbf{``Normality is needed for $b_0, b_1$ to be good.''} It is not:
    linearity, independence and constant variance suffice (Gauss--Markov).
    Normality gives the ML interpretation and exact small-sample tests.
  \item \textbf{Forgetting the logs}: maximise $\ln L$, not $L$ --- same
    maximiser, far easier algebra.
  \item \textbf{Unequal error variances}: if $\operatorname{Var}\varepsilon_i$
    differs by row, ML becomes \emph{weighted} least squares.
\end{itemize}
\end{caution}

\begin{tryit}[{{Exercises 19 --- numerical}}]
\begin{enumerate}[label=\textbf{E19.\arabic*}]
  \item Independent observations have densities $0.2$ and $0.3$ under a
    candidate model. Their joint likelihood, and its log?
  \item For the line of Example 18.3 ($x = (0,1,2,3)$, $y = (1,3,4,6)$,
    $\hat y = 1.1 + 1.6x$) compute SSE, $\hat\sigma^2_{\text{ML}}$, $s^2$ and $s$.
  \item For $x = (2,4,6,8,10)$, $y = (5,11,14,20,25)$ (E18.4) the direct slope is
    $2.45$. Compute the inverse slope $S_{yy}/S_{xy}$ and check that
    direct/inverse $= r^2$.
  \item Derive $\hat\sigma^2_{\text{ML}} = \mathrm{SSE}/n$ from $\ln L$.
\end{enumerate}
\end{tryit}

%======================================================================
\section{\texorpdfstring{$R^2$, adjusted $R^2$ and the error summary}{R squared, adjusted R squared and the error summary}}
\label{sec:r2}

\method{judging how good a fit is}
\begin{enumerate}
  \item $\mathrm{SST} = \sum(y_i - \bar y)^2$ (the ``predict $\bar y$''
    model), $\mathrm{SSE} = \sum e_i^2$, $\mathrm{SSR} = \sum(\hat y_i - \bar
    y)^2 = b_1^2S_{xx}$, and $\mathrm{SST} = \mathrm{SSR} + \mathrm{SSE}$
    (exact, because of the normal equations).
  \item $R^2 = \mathrm{SSR}/\mathrm{SST} = 1 - \mathrm{SSE}/\mathrm{SST}$: the
    share of the variation in $y$ the model accounts for. For one feature,
    $R^2 = r^2$ with $r = S_{xy}/\sqrt{S_{xx}S_{yy}}$ --- use it as a check.
  \item Adjusted: $R^2_{\text{adj}} = 1 - (1 - R^2)\frac{n-1}{n-p-1}$, $p$ =
    number of features (not counting the intercept). It can fall when a
    useless feature is added.
  \item On held-out data use the test mean in SST; $R^2$ can then be
    \emph{negative} (worse than predicting the mean).
  \item \textbf{Report three numbers}: $R^2$ (share explained), RMSE $=
    \sqrt{\mathrm{SSE}/n}$ (in the units of $y$, punishes big misses), MAE
    (typical miss).
\end{enumerate}
\end{keyidea}

\begin{worked}[{{Example 20.1 --- the six houses}}]
$\mathrm{SST} = S_{yy} = 2626$, $\mathrm{SSE} = 106$, $\mathrm{SSR} = 2520$
($= 0.06^2 \times 700000$ \checkmark).
\[
  R^2 = \tfrac{2520}{2626} = 0.9596, \qquad
  r = \tfrac{42000}{\sqrt{700000 \times 2626}} = 0.9796,\ r^2 = 0.9596\ \checkmark
\]
\[
  R^2_{\text{adj}} = 1 - (0.0404)\tfrac{5}{4} = 0.9495,
\]
\[
  \mathrm{RMSE} = \sqrt{106/6} = 4.20 \text{ lakh}, \qquad
  \mathrm{MAE} = \tfrac{22}{6} = 3.67 \text{ lakh}.
\]
``Explains $96\%$ of the variation'' and ``typically wrong by about 4 lakh'' are
both true; only the second tells a buyer whether the model is usable.
\end{worked}

\begin{worked}[{{Example 20.2 --- negative $R^2$, and adjusted $R^2$ deciding}}]
\textbf{Held-out rows} $y = (4, 6, 8, 10)$, predictions $(7, 6, 6, 5)$. Test
mean $7$: SST $= 9 + 1 + 1 + 9 = 20$; SSE $= 9 + 0 + 4 + 25 = 38$;
$R^2 = 1 - 38/20 = \mathbf{-0.90}$. Predicting $7$ everywhere would score $0$;
$R^2$ has no lower bound. (Likewise SSE $= 150$, SST $= 100$ gives
$R^2 = -0.5$.)

\textbf{Two models, $n = 30$}: A ($p = 2$, $R^2 = 0.64$) and B ($p = 8$,
$R^2 = 0.70$).
$R^2_{\text{adj}}(A) = 1 - 0.36 \times \tfrac{29}{27} = 0.6133$;
$R^2_{\text{adj}}(B) = 1 - 0.30 \times \tfrac{29}{21} = 0.5857$ --- A is
preferred. At $n = 200$ the same $R^2$s give $0.6363$ and $0.6874$ and B wins:
the penalty weakens as $n$ grows, so the final judge is still a held-out or
cross-validated score.
\end{worked}

\begin{caution}[{{Mistakes that cost marks}}]
\begin{itemize}
  \item \textbf{``$R^2 = 0.94$, so the model is 94\% accurate.''} $R^2$ is
    not accuracy (a classification term); it is unit-free and says nothing
    about the size of a typical error --- and a training $R^2$ is not evidence
    at all.
  \item \textbf{Training $R^2$ rose when I added a feature, so it helps.} It
    can never fall: the old fit is still available with the new coefficient at
    $0$.
  \item \textbf{Test $R^2$ with the training mean}, or $p$ counting the
    intercept.
  \item \textbf{$R^2$ as proof the form is right.} Anscombe's four datasets
    share $R^2 \approx 0.67$; only one is a straight line with scatter.
\end{itemize}
\end{caution}

\begin{tryit}[{{Exercises 20 --- numerical}}]
\begin{enumerate}[label=\textbf{E20.\arabic*}]
  \item SSE $= 20$ and SST $= 100$ on the same evaluation set: $R^2$? And for
    SSE $= 150$, SST $= 100$?
  \item For the fit of E18.4 ($\hat y = 0.30 + 2.45x$): SSE, SST, $R^2$ (two
    ways), RMSE and MAE.
  \item $n = 50$, $p = 6$, SST $= 48{,}000$, SSE $= 7{,}200$: $R^2$, adjusted
    $R^2$ and $s = \sqrt{\mathrm{SSE}/(n - p - 1)}$. A seventh feature raises
    $R^2$ to $0.858$: the new adjusted $R^2$ --- did it earn its place?
  \item Held-out $y = (2, 5, 8)$, predictions $(6, 5, 3)$: the test $R^2$.
  \item $n = 20$. A: $p = 1$, $R^2 = 0.50$. B: $p = 6$, $R^2 = 0.60$. Which
    does adjusted $R^2$ prefer?
\end{enumerate}
\end{tryit}

%======================================================================
\section{Model adequacy, over-fitting and cross-validation}
\label{sec:adequacy}

\method{checking a fitted regression}
\begin{enumerate}
  \item \textbf{Plot residuals against fitted values.} A formless band around
    $0$ is healthy. A curve (smile or frown) $\Rightarrow$ wrong functional
    form (add $x^2$, transform). A funnel $\Rightarrow$ non-constant variance
    (log the target, or weighted least squares). Isolated points $\Rightarrow$
    investigate them; a high-\emph{leverage} point (extreme in $x$) can rotate
    the line.
  \item Residuals against each feature, against time (drift $\Rightarrow$
    dependence), and a normal Q--Q plot (least important).
  \item \textbf{Over-fitting}: great training score, much worse held-out score.
    \textbf{Under-fitting}: poor on both. Over-fitting depends on flexibility
    \emph{relative to} $n$.
  \item \textbf{Cross-validation as the detector}: choose the model (e.g.\ the
    polynomial degree) with the best cross-validated error; report the mean and
    spread across folds. Pool fold errors through MSE: pooled RMSE $=
    \sqrt{\text{mean of fold MSEs}}$ (equal-size folds).
\end{enumerate}
\end{keyidea}

\begin{worked}[{{Example 21.1 --- the degree curve}}]
Sixteen houses fitted, eight withheld; polynomials in area:
\begin{center}
\small
\begin{tabular}{@{}rrrl@{}}
\toprule
degree & train $R^2$ & held-out $R^2$ & verdict\\ \midrule
1 & $0.811$ & $+0.838$ & the relationship, correctly identified\\
3 & $0.820$ & $+0.758$ & extra flexibility already costing accuracy\\
5 & $0.954$ & $-37.5$ & fitting noise\\
9 & $0.979$ & $-16{,}028$ & fitting noise energetically\\
\bottomrule
\end{tabular}
\end{center}
Training $R^2$ climbs every time; the held-out score collapses. Every training-set
score ($R^2$, even adjusted $R^2$) prefers degree 9 --- over-fitting is
invisible from inside. The fix is procedural: withhold data and never let a
decision depend on it. Choose \textbf{degree 1}.
\end{worked}

\begin{worked}[{{Example 21.2 --- reading fold scores}}]
Five folds give RMSE $3.0, 2.5, 3.5, 4.0, 3.0$ lakh. Mean fold RMSE $= 3.2$.
Pooled RMSE $= \sqrt{(9 + 6.25 + 12.25 + 16 + 9)/5} = \sqrt{10.5} = 3.2404$
--- not the same number, because RMSE is a square root of an average. The
spread ($2.5$ to $4.0$) is part of the answer: quoting one split's RMSE to four
figures claims a precision the data does not have.

\textbf{Model A}: train RMSE $2$, validation RMSE $9$. \textbf{Model B}: $4$
and $5$. Choose \textbf{B}; A's training error flatters it (over-fitting).
\end{worked}

\begin{tryit}[{{Exercises 21 --- numerical}}]
\begin{enumerate}[label=\textbf{E21.\arabic*}]
  \item Fold RMSEs are $2, 3, 4$: their mean? Three equal folds have MSEs $4, 9,
    16$: the pooled out-of-fold RMSE?
  \item Three models: A (train $R^2$ $0.63$, test $0.61$), B ($0.88$, $0.85$),
    C ($0.999$, $0.41$). Diagnose each and choose.
\end{enumerate}
\end{tryit}

\begin{concept}
\begin{enumerate}[label=\textbf{C21.\arabic*}]
  \item \emph{A curved residual-vs-fitted pattern suggests?} The linear form
    is wrong (linearity violated); the coefficients describe a relationship
    that does not exist. Add a polynomial term or transform the feature.
  \item \emph{A funnel?} Non-constant variance: coefficients stay unbiased but
    their standard errors and intervals are wrong. Log the target or use
    weighted least squares.
  \item \emph{Distinguish nonlinearity, unequal variance and dependence in
    diagnostics. Can CV alone establish every assumption?} Curve vs funnel vs
    pattern over time/order. No --- CV measures predictive error; it does not
    test independence or constant variance.
  \item \emph{Why is $R^2$ undefined when all evaluated targets are equal? Why
    can training $R^2$ rise with a useless predictor?} SST $= 0$, so the ratio
    divides by zero. Adding a column enlarges the space least squares
    searches, so SSE cannot rise.
  \item \emph{Adding a feature raises training $R^2$ from $0.81$ to $0.83$
    but lowers CV $R^2$ from $0.79$ to $0.74$.} Over-fitting: drop the
    feature; trust the cross-validated score.
  \item \emph{Name Anscombe's point in one sentence.} Four datasets with the
    same means, line and $R^2$ look completely different when plotted ---
    always plot.
  \item \emph{Leave-one-out or 5-fold?} 5- or 10-fold is the default; LOOCV
    ($n$ fits) is nearly unbiased but high-variance, for very small datasets.
\end{enumerate}
\end{concept}
%======================================================================
\section{Practice papers}
\label{sec:papers}

Sit each paper in 60 minutes without notes, then mark it against the
solution. Both follow the class-test pattern: five questions, 20 marks,
simple calculator allowed, all steps shown.

\subsection*{Class Test --- Set B \quad (as set on 1 September 2026)}

\begin{enumerate}[label=\textbf{Q\arabic*.}]
  \item \textbf{Standardisation} \hfill [3]\\
    In a housing dataset the Price column has mean $80$ lakh and standard
    deviation $20$ lakh.
    (a) Give one difference between normalisation (min--max) and
    standardisation ($z$-score). (1)
    (b) Standardise three houses priced $100$, $60$ and $80$ lakh. (2)
  \item \textbf{Data reduction and sampling} \hfill [3]\\
    A dataset has 400 houses and 15 columns, one of which is House\_ID.
    (a) Why should House\_ID be removed before training? Which preprocessing
    step is this? (2)
    (b) Only 40 of the 400 houses are marked expensive. Name one sampling
    method to fix this. Does it add rows or remove rows? (1)
  \item \textbf{Model validation} \hfill [4]\\
    (a) Give the use of the training, validation and test sets. (2)
    (b) A student says ``my model is 96\% accurate'', measured on the
    training data. Why is this number not correct, and what should he do
    instead? (1)
    (c) 400 houses are split $75{:}25$. How many go in each part? (1)
  \item \textbf{Finding the formulas for $b_0$ and $b_1$} \hfill [4]\\
    We fit $\hat y = b_0 + b_1x$. Write down what is made as small as
    possible, derive $b_0$ and $b_1$ showing all steps, and show that the line
    passes through $(\bar x, \bar y)$.
  \item \textbf{Fitting a straight line} \hfill [6]\\
    \begin{tabular}{@{}lccccc@{}}
    Area $x$ (sq ft) & 600 & 900 & 1200 & 1500 & 1800\\
    Price $y$ (lakh) & 42 & 65 & 80 & 97 & 116
    \end{tabular}\\
    (a) Find $\bar x$, $\bar y$, $S_{xx}$ and $S_{xy}$. (3)
    (b) Find $b_1$ and $b_0$ and write the line. (2)
    (c) Predict the price of a 1350 sq ft house. (1)
\end{enumerate}

\begin{worked}[{{Set B --- solutions}}]
\textbf{Q1.} (a) Min--max maps to a fixed range $[0,1]$ using the minimum and
maximum; standardisation centres on the mean and divides by the standard
deviation, giving mean $0$, sd $1$ and no fixed range (it is less distorted by
a single extreme than min--max). (b) $z = (x - 80)/20$: $100 \to +1$;
$60 \to -1$; $80 \to 0$.

\textbf{Q2.} (a) An identifier carries no information about price; a model can
still latch onto it (memorising rows, or spurious order), which hurts
generalisation. Removing it is \emph{data reduction --- feature selection}
(dropping an irrelevant column). (b) Random oversampling (or SMOTE) of the
expensive class --- \emph{adds} rows; or random undersampling of the rest ---
\emph{removes} rows.

\textbf{Q3.} (a) Training: learn the parameters. Validation: compare models
and tune hyperparameters. Test: one final, unbiased estimate on unseen data.
(b) A training score measures how well the model fits data it has already
seen and rewards memorisation (overfitting); evaluate on a held-out test set
or by cross-validation. (c) $300$ training, $100$ test.

\textbf{Q4.} See Example 18.1: minimise $S = \sum(y_i - b_0 - b_1x_i)^2$;
normal equations; $b_0 = \bar y - b_1\bar x$, $b_1 = S_{xy}/S_{xx}$; at
$x = \bar x$, $\hat y = \bar y$.

\textbf{Q5.} (a) $\bar x = 6000/5 = 1200$, $\bar y = 400/5 = 80$.
$x - \bar x = (-600, -300, 0, 300, 600)$; $y - \bar y = (-38, -15, 0, 17,
36)$. $S_{xx} = 360000 + 90000 + 0 + 90000 + 360000 = 900000$;
$S_{xy} = 22800 + 4500 + 0 + 5100 + 21600 = 54000$.
(b) $b_1 = 54000/900000 = 0.06$; $b_0 = 80 - 0.06 \times 1200 = 8$;
$\hat y = 8 + 0.06x$. (Check: $\hat y = 44, 62, 80, 98, 116$; residuals
$-2, 3, 0, -1, 0$, sum $0$.) (c) $\hat y(1350) = 8 + 81 = 89$ lakh.
\end{worked}

\subsection*{Practice Set C \quad (new, same pattern)}

\begin{enumerate}[label=\textbf{Q\arabic*.}]
  \item \textbf{Scaling} \hfill [3]\\
    Training Area has mean $1500$ and sd $400$ sq ft (population); its minimum
    is $1000$ and maximum $3000$.
    (a) Standardise areas $1900$, $1100$ and $1500$. (1.5)
    (b) A test house has area $3400$. Give its min--max scaled value, and say
    why the answer exceeding $1$ is not an error. (1.5)
  \item \textbf{Imbalanced data} \hfill [4]\\
    600 loan applications, 60 of which defaulted.
    (a) What accuracy does ``nobody defaults'' achieve? What is its recall? (1)
    (b) A model gives $TP = 45$, $FN = 15$, $FP = 55$, $TN = 485$. Compute
    accuracy, precision, recall and $F_1$. (2)
    (c) How many rows after random oversampling to balance? After random
    undersampling? Where in the workflow must resampling happen? (1)
  \item \textbf{Optimisers} \hfill [4]\\
    $J(w) = \tfrac12(w - 5)^2$, $w_0 = 1$, $\eta = 0.5$.
    (a) Two steps of plain gradient descent. (1.5)
    (b) Two steps of momentum with $\beta = 0.5$ ($v \leftarrow \beta v + g$). (1.5)
    (c) The first Adam step with default $\beta$'s (ignore $\varepsilon$). (1)
  \item \textbf{Missing values and outliers} \hfill [3]\\
    A column reads $30, 34, ?, 38, 40, 41, 44, 46, 120$.
    (a) Mean and median of the observed values; which would you impute, and
    why? (1)
    (b) Quartiles (hinges), IQR and fences; which values are flagged? (1.5)
    (c) Name one way to keep the flagged value but reduce its influence. (0.5)
  \item \textbf{Fitting a straight line} \hfill [6]\\
    \begin{tabular}{@{}lccccc@{}}
    Experience $x$ (years) & 1 & 3 & 5 & 7 & 9\\
    Salary $y$ (thousand/month) & 29 & 37 & 42 & 49 & 53
    \end{tabular}\\
    (a) $\bar x$, $\bar y$, $S_{xx}$, $S_{xy}$. (2)
    (b) $b_1$, $b_0$, the line, and a check on the residuals. (1.5)
    (c) Predict the salary at 6 years. (1)
    (d) Compute SSE, $R^2$ and $s = \sqrt{\mathrm{SSE}/(n-2)}$. (1.5)
\end{enumerate}

\begin{worked}[{{Set C --- solutions}}]
\textbf{Q1.} (a) $z = (x - 1500)/400$: $1900 \to 1$, $1100 \to -1$,
$1500 \to 0$. (b) $(3400 - 1000)/2000 = 1.2$. The scaler is fitted on
training data only, so a test value beyond the training maximum maps above
$1$; that correctly signals extrapolation. Refitting on the test set would be
leakage.

\textbf{Q2.} (a) $540/600 = 90\%$; recall $0$. (b) Accuracy $530/600 =
0.8833$; precision $45/100 = 0.45$; recall $45/60 = 0.75$; $F_1 =
2(0.45)(0.75)/1.20 = 0.5625$. Accuracy is \emph{lower} than the useless
model's, yet it finds $75\%$ of defaulters. (c) Oversampling: $540 + 540 =
1080$; undersampling: $60 + 60 = 120$. Only on the training data, after the
split (inside each CV fold).

\textbf{Q3.} $J'(w) = w - 5$. (a) $g = -4 \Rightarrow w = 1 + 2 = 3$;
$g = -2 \Rightarrow w = 4$. (b) $v_1 = -4$, $w_1 = 3$; $g_2 = -2$,
$v_2 = 0.5(-4) - 2 = -4$, $w_2 = 3 + 2 = 5$ --- momentum reaches the minimum
in two steps where GD is at $4$. (c) $\hat m_1 = g_1 = -4$, $\hat v_1 = 16$,
step $= -0.5(-4)/4 = +0.5$, $w_1 = 1.5$ --- exactly $\eta$.

\textbf{Q4.} (a) Observed: $30, 34, 38, 40, 41, 44, 46, 120$; mean
$393/8 = 49.125$, median $(40 + 41)/2 = 40.5$. The $120$ drags the mean above
seven of eight values: impute the \textbf{median}. (b) Lower half $30, 34, 38,
40 \Rightarrow Q_1 = 36$; upper half $41, 44, 46, 120 \Rightarrow Q_3 = 45$;
$\IQR = 9$; fences $22.5$ and $58.5$: $120$ is flagged. (c) Clip it to $58.5$,
log-transform the column, use a robust scaler, or use a robust loss/model.

\textbf{Q5.} (a) $\bar x = 5$, $\bar y = 42$; $x - \bar x = (-4,-2,0,2,4)$,
$y - \bar y = (-13,-5,0,7,11)$; $S_{xx} = 40$, $S_{xy} = 52 + 10 + 0 + 14 + 44
= 120$. (b) $b_1 = 3$, $b_0 = 42 - 15 = 27$; $\hat y = 27 + 3x$; fitted
$30, 36, 42, 48, 54$, residuals $-1, 1, 0, 1, -1$, sum $0$ \checkmark.
(c) $\hat y(6) = 45$ thousand per month.
(d) SSE $= 1 + 1 + 0 + 1 + 1 = 4$; SST $= S_{yy} = 169 + 25 + 0 + 49 + 121 =
364$; $R^2 = 1 - 4/364 = 0.9890$ (check: $r = 120/\sqrt{40 \times 364} =
0.9945$, $r^2 = 0.9890$); $s = \sqrt{4/3} = 1.1547$ thousand.
\end{worked}
%======================================================================
\section{Answers to the numerical exercises}
\label{sec:answers}

Rounded to four decimals at the end. Compare the \emph{method} line by line,
not only the final value. (The short-answer questions have their model
answers beside them in the grey boxes.)

\ans{Exercises 1}
\textbf{E1.1} $5 + 0.06 \times 1000 = 65$ lakh.
\textbf{E1.2} $120 \times 4 = 480$ entries; $X$ is $120 \times 4$, $y$ has
$120$ entries.
\textbf{E1.3} Distances $0.1414, 0.2236, 3.3838, 3.2311, 5.0090$: nearest is
$(1.4, 0.2)$ --- \textbf{setosa}.
\textbf{E1.4} $94\%$ are not readmitted, so predicting ``not readmitted'' for
almost everyone scores about $94\%$ while missing most real cases. Recall on the
readmitted class (here $20\%$), or the confusion matrix, exposes it.

\ans{Exercises 2}
\textbf{E2.1} $50 - 47 = 3$. $40 - 45 = -5$: over-predicting by $5$.
\textbf{E2.2} MSE $= (1 + 4 + 16)/3 = 7$; MAE $= 7/3 = 2.3333$; Huber
$0.5, 1.5, 3.5$, mean $5.5/3 = 1.8333$.
\textbf{E2.3} $|1| \le 2$: $0.5$; $|-3| > 2$: $2(3 - 1) = 4$; mean $2.25$.
\textbf{E2.4} MAE $= 3$, MSE $= (16 + 4)/2 = 10$. Squared penalty $9 \to 36$
($\times 4$); absolute $3 \to 6$ ($\times 2$).
\textbf{E2.5} $e = (-1, 0, 2, -4)$: MSE $= 21/4 = 5.25$, MAE $= 1.75$,
RMSE $= 2.2913$; Huber $0.5, 0, 2, 6$, mean $2.125$. Row 4 supplies
$16/21 = 76.2\%$ of the MSE.
\textbf{E2.6} MSE: mean $= 4$; MAE: median $= 2$. With $112$: mean $24$
(moved by $20$), median still $2$.

\ans{Exercises 3}
\textbf{E3.1} $J'(5) = 2$; $w = 5 - 0.2 = 4.8$.
\textbf{E3.2} $w_{\text{new}} - 3 = (1 - \eta)(w_{\text{old}} - 3)$:
converges for $0 < \eta < 2$; one step at $\eta = 1$.
\textbf{E3.3} $J'(1) = -1$: subtracting a negative gradient moves right,
towards $2$. From $5$: $g = 3$, $w = 4.7$.
\textbf{E3.4} Factor $1 - 2(0.2) = 0.6$. $w = 3, 1.4, 0.44, -0.136, -0.4816,
-0.6890$; $f = 16, 5.76, 2.0736, 0.7465, 0.2687, 0.0967$.
\textbf{E3.5} $J_0 = 11.6667$; gradient $(-6, -8.6667) \Rightarrow (b, w) =
(0.6, 0.8667)$, $J = 3.2074$; gradient $(-3.0667, -4.5778) \Rightarrow
(0.9067, 1.3244)$, $J = 0.8954$.

\ans{Exercises 4}
\textbf{E4.1} $5$.
\textbf{E4.2} Row gradients $3$ and $1$; mean $2$.
\textbf{E4.3} Full gradient $(1 + (-1))/2 = 0$; SGD gradient for $y = 1$ is
$+1$. Yes: with $\eta = 0.2$, $w \to 1.8$ and the full loss rises $0.5 \to
0.52$.
\textbf{E4.4} $g = 3$, $w = 6 - 0.6 = 5.4$.
\textbf{E4.5} $w^\star = 26/14 = 1.8571$. Batch: mean of $(-16, -6, -30)$ is
$-17.3333$, $w = 0.8667$. SGD: $0.8, 1.02, 1.602$ --- much nearer.
\textbf{E4.6} $200$ updates per epoch, $2000$ in all; $500{,}000$
row-gradients.

\ans{Exercises 5}
\textbf{E5.1} $v_1 = 2$, $v_2 = 0.9(2) + 2 = 3.8$.
\textbf{E5.2} $v_t = 2(1 + 0.9 + \dots + 0.9^{t-1}) = 20(1 - 0.9^t) \to 20$.
\textbf{E5.3} $g = -6, v = -6, w = 0.6$; $g = -4.8, v = -7.8, w = 1.38$;
$g = -3.24, v = -7.14, w = 2.094$ (GD: $1.464$).
\textbf{E5.4} $4,\ 6.2,\ 3.96,\ 1.168$.

\ans{Exercises 6}
\textbf{E6.1} $9 + 16 = 25$; $w = 2 - 0.1 \times 4/5 = 1.92$.
\textbf{E6.2} $V = 4$, $w = 5 - 0.1 \times 2/2 = 4.9$. From $V = 0$:
$V = 0.4$, $\sqrt V = 0.6325$ (not $\sqrt4$).
\textbf{E6.3} $g = 8, G = 64, w = 4.5$; $g = 7, G = 113, \sqrt G = 10.6301,
w = 4.1707$; $g = 6.3415, G = 153.2146, \sqrt G = 12.3780, w = 3.9146$.
\textbf{E6.4} Adagrad $G = 16, 20, 21, 21.25$; ratios $1, 0.4472, 0.2182,
0.1085$. RMSprop $V = 1.6, 1.84, 1.756, 1.6054$; ratios $3.1623, 1.4744,
0.7546, 0.3946$. RMSprop's stays larger: it forgets the early large gradients.
\textbf{E6.5} $V = \rho^kV_0$. No: the numerator $g = 0$, so $w$ does not move
--- only the history changes.

\ans{Exercises 7}
\textbf{E7.1} $m_1 = 0.2$, $v_1 = 0.004$, $\hat m_1 = 2$, $\hat v_1 = 4$,
$w_1 = 4.9$. With $g_1 = -2$: $\hat m_1 = -2$, $\hat v_1 = 4$, step $+0.1$.
\textbf{E7.2} $g = 4$: $m = 0.4$, $v = 0.016$, $\hat m = 4$, $\hat v = 16$,
$w = 2.9$. $g = 3.8$: $m = 0.74$, $v = 0.030424$, $\hat m = 3.8947$,
$\hat v = 15.2196$, $\sqrt{\hat v} = 3.9012$, $w = 2.8002$.
\textbf{E7.3} $\hat m_1 = (1-\beta_1)g_1/(1-\beta_1) = g_1$,
$\hat v_1 = g_1^2$, step $= -\eta g_1/|g_1|$. For $g_1 = -0.003$: $+0.01$.
\textbf{E7.4} $(1-\beta_1)/\sqrt{1-\beta_2}$: $0.1/0.0316 = 3.1623\eta$;
with $\beta_2 = 0.99$, $0.1/0.1 = 1\eta$.

\ans{Exercises 8}
\textbf{E8.1} $G = 4$, ratio $1/2$, $\Delta w = -1$.
\textbf{E8.2} $G = 0.9 + 0.4 = 1.3$, $\sqrt G = 1.1402$;
$\Delta w = -(0.5/1.1402)(-2) = +0.8771$; $U = 0.9(0.25) + 0.1(0.7692) =
0.3019$.

\ans{Exercises 9}
\textbf{E9.1} Median of $10, 20, 40$ is $20$.
\textbf{E9.2} Mean $18000/6 = 3000$; median $(1800 + 2000)/2 = 1900$; impute
the median --- the 9000 makes the mean unrepresentative.
\textbf{E9.3} Mean $299.5/11 = 27.23$, median $18$; the $95$ skews the mean,
so fill with $18$ (add a missing-indicator if missingness may be informative).
H12 is a modest Prem Nagar house, so $18$ is also plausible --- though Income
is a leak and should be dropped anyway.
\textbf{E9.4} Population sd $\times\sqrt{80/100} = 0.8944$; sample sd
$\times\sqrt{79/99} = 0.8933$.

\ans{Exercises 10}
\textbf{E10.1} $\IQR = 8$; upper fence $18 + 12 = 30$.
\textbf{E10.2} $Q_1 = (25 + 30)/2 = 27.5$, $Q_3 = (45 + 50)/2 = 47.5$,
$\IQR = 20$, fences $-2.5$ and $77.5$: $200$ flagged.
\textbf{E10.3} Halves $12, 14, 15, 15$ and $17, 18, 19, 60$: $Q_1 = 14.5$,
$Q_3 = 18.5$, $\IQR = 4$, fences $8.5$ and $24.5$: flagged.
$\bar x = 20.6667$, $s = 14.8997$, $z = 2.6399 < 3$: not flagged --- and never
could be, since the ceiling for $n = 9$ is $8/3 = 2.6667$; the $60$ inflates
the $\bar x$ and $s$ it is judged by. Median $16$, $\MAD = 2$, $M = 0.6745
\times 44/2 = 14.84$: flagged.
\textbf{E10.4} $9/\sqrt{10} = 2.8460$; $10/\sqrt{11} = 3.0151$. Only from
$n = 11$ upwards.
\textbf{E10.5} $\bar x = 2$, $m_2 = 16$, $m_3 = 96$, $g_1 = 96/64 = 1.5$ ---
exactly the ceiling $3/\sqrt4$: one point away from four equal ones is as
skewed as five points can be.

\ans{Exercises 11}
\textbf{E11.1} $(100 - 80)/20 = 1$; $(40 - 10)/20 = 1.5$.
\textbf{E11.2} $1500/4500 = 0.3333$.
\textbf{E11.3} $\mu = 30$, $\sigma = \sqrt{1000/5} = 14.142$,
$z_{50} = 1.4142$.
\textbf{E11.4} $(60 - 37.5)/20 = 1.125$.
\textbf{E11.5} Unscaled: $d_A = 50.04$, $d_B = 400$ --- A. Scaled: Q $(0.4,
0.5)$, A $(0.425, 0)$, B $(0.2, 0.5)$; $d_A = 0.5006$, $d_B = 0.2$ --- B.
\textbf{E11.6} Mean $4.6667$, $\sigma = 6.3944$; standard $-1.3553, -0.7298,
-0.4170, 0.2085, 0.5213, 1.7724$. Min--max $(v + 4)/20$: $0, 0.2, 0.3, 0.5,
0.6, 1$. Median $4$, $Q_1 = 0$, $Q_3 = 8$, robust $(v - 4)/8$: $-1, -0.5,
-0.25, 0.25, 0.5, 1.5$. Max-abs $v/16$: $-0.25, 0, 0.125, 0.375, 0.5, 1$.

\ans{Exercises 12}
\textbf{E12.1} $3$.
\textbf{E12.2} $6 - 1 + 5 = 10$; with drop-first $9$.
\textbf{E12.3} $\bar y = 20$, $S_{xy} = -10$, $S_{xx} = 2$: $\hat y = 25 - 5x$,
fitted $25, 20, 15$; SSE $= 25 + 100 + 25 = 150$ (predicting the mean $20$
everywhere gives $200$). One-hot fits each mean exactly: SSE $= 0$.

\ans{Exercises 13}
\textbf{E13.1} $(2 + 2)/\sqrt2 = 2.8284$.
\textbf{E13.2} $\begin{pmatrix}1 & 0.8\\0.8 & 1\end{pmatrix}$; $\lambda = 1.8,
0.2$; PC1 explains $90\%$.
\textbf{E13.3} Mean $(5, 4)$; $S = \begin{pmatrix}6 & 2\\2 & 3\end{pmatrix}$;
$\lambda^2 - 9\lambda + 14 = 0 \Rightarrow 7, 2$; $u_1 = (2,1)/\sqrt5$,
$u_2 = (-1, 2)/\sqrt5$; $77.8\%$; scores $(-9,-1,0,3,7)/\sqrt5 =
(-4.0249, -0.4472, 0, 1.3416, 3.1305)$, variance $7$.
\textbf{E13.4} Cumulative $0.52, 0.73, 0.84, 0.93, 0.97$: $4$ for $90\%$, $5$
for $95\%$.

\ans{Exercises 14}
\textbf{E14.1} $300$.
\textbf{E14.2} $20$ rows per fold; $5$ fits.
\textbf{E14.3} Mean $0.80$; deviations $0.01, -0.01, 0.06, -0.06, 0$; sd
$= \sqrt{0.0074/5} = 0.0385$.
\textbf{E14.4} $\binom{24}{5}/\binom{30}{5} = 42504/142506 = 0.2983$;
stratified: $1$ positive and $4$ negatives.

\ans{Exercises 15}
\textbf{E15.1} $10\%$; accuracy $90\%$.
\textbf{E15.2} $TP = 30$, $FP = 60$, $FN = 10$, $TN = 1900$; accuracy $0.965$,
precision $0.3333$, recall $0.75$, $F_1 = 0.4615$.
\textbf{E15.3} $TP = 60$, $FP = 90$, $FN = 40$, $TN = 9810$; accuracy $0.987$,
precision $0.40$, recall $0.60$, $F_1 = 0.48$. The rival: accuracy $99\%$,
recall $0$ --- deploy the first.
\textbf{E15.4} $2400$ rows; $160$ rows, discarding $1120$.
\textbf{E15.5} Accuracy $98\%$; recall $0$.

\ans{Exercises 18}
\textbf{E18.1} $4$; $106/6 = 17.67$.
\textbf{E18.2} $\bar x = 2$, $\bar y = 3.3333$, $S_{xx} = 2$, $S_{xy} = 3$:
$b_1 = 1.5$, $b_0 = 0.3333$.
\textbf{E18.3} $40 - 45 = -5$.
\textbf{E18.4} $\bar x = 6$, $\bar y = 15$, $S_{xx} = 40$, $S_{xy} = 98$,
$b_1 = 2.45$, $b_0 = 0.30$; $\hat y(7) = 17.45$. Residuals $-0.2, 0.9, -1.0,
0.1, 0.2$ (sum $0$); SSE $= 1.90$.
\textbf{E18.5} $S_{xx} = 310 - 250 = 60$, $S_{xy} = 660 - 600 = 60$, $b_1 = 1$,
$b_0 = 12 - 5 = 7$, $\hat y(7) = 14$.
\textbf{E18.6} $w = 35/14 = 2.5$; residuals $1.5, 0, -0.5$, sum $1 \ne 0$.
$\sum xe = 0$ survives (the normal equation for $w$); ``sum to zero'' is the
intercept's normal equation, and there is no intercept.

\ans{Exercises 19}
\textbf{E19.1} $0.2 \times 0.3 = 0.06$; $\ln 0.06 = -2.8134$.
\textbf{E19.2} Residuals $-0.1, 0.3, -0.3, 0.1$; SSE $= 0.2$;
$\hat\sigma^2_{\text{ML}} = 0.05$; $s^2 = 0.2/2 = 0.1$; $s = 0.3162$.
\textbf{E19.3} $S_{yy} = 242$, $S_{xy} = 98$: inverse slope $242/98 = 2.4694$;
$2.45/2.4694 = 0.9921 = 98^2/(40 \times 242) = r^2$ \checkmark.
\textbf{E19.4} $\partial\ln L/\partial\sigma^2 = -\frac{n}{2\sigma^2} +
\frac{\mathrm{SSE}}{2\sigma^4} = 0 \Rightarrow \sigma^2 = \mathrm{SSE}/n$.

\ans{Exercises 20}
\textbf{E20.1} $1 - 20/100 = 0.8$; $1 - 150/100 = -0.5$ (worse than predicting
the mean).
\textbf{E20.2} SSE $= 1.90$; SST $= 100 + 16 + 1 + 25 + 100 = 242$;
$R^2 = 1 - 1.90/242 = 0.9921$ ($r = 98/\sqrt{9680} = 0.9961$, $r^2 = 0.9921$);
RMSE $= \sqrt{1.90/5} = 0.6164$; MAE $= 2.40/5 = 0.48$.
\textbf{E20.3} $R^2 = 0.85$; $R^2_{\text{adj}} = 1 - 0.15 \times 49/43 = 0.8291$;
$s = \sqrt{7200/43} = 12.94$. With $p = 7$: $1 - 0.142 \times 49/42 = 0.8343$
--- it rose, so on this criterion the feature earned its place (confirm on
held-out data).
\textbf{E20.4} Test mean $5$: SST $= 18$, SSE $= 16 + 0 + 25 = 41$,
$R^2 = -1.2778$.
\textbf{E20.5} A: $1 - 0.5 \times 19/18 = 0.4722$; B: $1 - 0.4 \times 19/13 =
0.4154$. A.

\ans{Exercises 21}
\textbf{E21.1} Mean $3$; pooled $\sqrt{(4 + 9 + 16)/3} = \sqrt{29/3} = 3.1091$.
\textbf{E21.2} A under-fits (both low, small gap); B fits well; C over-fits
(near-perfect training, collapsed test). Deploy B; for A add features or
flexibility.

\vspace{1em}
\begin{keyidea}[{{The night before}}]
\begin{itemize}
  \item Redo one worked example per section \emph{without looking}: a Huber
    table, a three-step Adam table, the IQR fences, a KNN distance before and
    after scaling, a $2\times2$ PCA, a confusion matrix, the six-house fit.
  \item Rehearse the derivation of $b_0$ and $b_1$, and the step from the
    Gaussian log-likelihood to least squares, until you can write them
    unprompted --- both are favourite 4-mark questions.
  \item For any fitted line, be able to produce SSE, SST, SSR, $R^2$,
    adjusted $R^2$, $\hat\sigma^2_{\text{ML}}$, $s$, RMSE and MAE in one pass.
  \item Learn the one-line distinctions: selection cuts columns, sampling cuts
    rows; normalise to a range, standardise to a mean; an error is wrong, an
    outlier is rare; validation chooses, test reports once; resample the
    training set only.
  \item Out of syllabus for this paper: cross-entropy, hinge and triplet
    losses, and everything from Lecture 10 on (logistic regression, multiple
    regression, regularisation).
\end{itemize}
\end{keyidea}

\end{document}
